\documentclass[a4paper]{report}
\usepackage{amsmath,amssymb}
\usepackage{graphicx}
\graphicspath{ {./picture/} }
\usepackage{fancyhdr}
%\usepackage[colorlinks,urlcolor=blue]{hyperref}
\usepackage{url}
\usepackage{tikz}
\usetikzlibrary{arrows}
\usetikzlibrary{positioning} % This line includes the positioning library

%%%%%% code 
\usepackage{listings}


% این فایل شامل تنظیماتی است که قبل از لود شدن پکیج زی‌پرشین باید انجام شوند.
% دلیل آن است که نیاز به لود شدن یک پکیج (geometry) دارد، در حالی که bidi (که توسط زی‌پرشین
% لود می‌شود) حتما باید آخرین پکیج باشد. در نتیجه هر پکیج دیگری نیز که بخواهد لود شود، باید در این قسمت باشد.
% نسخه 1.0.7
% تاریخ: ۲۲ مهر ۱۳۹۳
% تغییرات: ۱) برای کارکردن با texlive2014 و زی‌پرشین جدید:
%               - فونتها با IRLotus جایگزین شدند => فقط با texlive2013 و texlive2014 کار می‌کند (با 2012
%                  کار نمی‌کند، چون کشیده و کاراکتر نیمفاصله را در خروجی بین دو عدد "" می‌گذارد!).
%               - بجای داشتن یک فایل MBZ_ThesisSettings.sty، دو فایل MBZ_ThesisSettings1.tex و
%                  MBZ_ThesisSettings1.tex داریم، که اولی شامل تنظیماتی است که قبل از لود کردن زی‌پرشین
%                  باید انجام شوند و دومی تنظیماتی که بعد از آن باید انجام شوند.

%%%%%%%%%%%%%%%%%%%%%%%%%%%%%%%%%%%%%%%%%%
% Start of Page Setup:
% تعریف متغیرها:
\newlength{\MyLeftMargin}
\newlength{\MyRightMargin}
\newlength{\MyTopMargin}
\newlength{\MyBottomMargin}

% اعداد در استاندارد تزهای دانشگاه شریف:
\setlength{\MyLeftMargin}{2.5cm}
\setlength{\MyRightMargin}{3cm} % -> Sharif says 3cm. I prefer 2.5.
\setlength{\MyTopMargin}{3cm}
\setlength{\MyBottomMargin}{2.5cm} % -> Sharif says 2.5cm. I prefer 2.7.

\usepackage[top=\MyTopMargin, bottom=\MyBottomMargin, left=\MyLeftMargin, right=\MyRightMargin]{geometry}
%\setlength{\columnwidth}{82mm}
%\setlength{\columnsep}{6mm}
% End of Page Setup
%%%%%%%%%%%%%%%%%%%%%%%%%%%%%%%%%%%%%%%%%% % All settings before xepersian (especially loading any other pachages,
                            % because xepersian should be the last loaded package.
							% written by Massoud Babaie-Zadeh
\usepackage[Kashida]{xepersian}
% این فایل شامل تنظیمات نوشته‌شده توسط اینجانب (مسعود بابایی‌زاده یا  MBZ) برای نوشتن تز یا گزارش است، چون خروجی پیش‌فرض xepersian اصلا تز زیبایی تولید نمی‌کند. همچنین دقت شود که از فونت MBZ Lotus و نیز پکیج digitfont (برای تغییر فونت ارقام به MBZ Lotus) که هر دو توسط اینجانب نوشته شدند به همراه این تنظیمات استفاده شود تا در نهایت خروجی زیبایی برای تز بدست آید.
% تنظیمات انجام شده در این فایل شامل موارد زیر هستند:
% اول: زیاد کردن فاصله خطوط در متن فارسی، کمتر زیاد کردن در ریاضی و متن انگلیسی: چون در XePersian زیادی خطوط توی هم هستند و زیبا نیست. این کار کار راحتی نبود و خیلی چیزها دست خوردند.
% دوم: فهرست مطالب در xepersian به نظرم اصلا زیبا نیست، چون در واقع چیز خاصی برای فارسی ندارد و از همان report استفاده می‌کند که برای زیبایی در انگلیسی طراحی شده نه فارسی. در نتیجه تغییرات متعددی در آن دادم تا فهرست مطالب به صورت زیباتری نوشته شود.
% سوم: تنظیم صفحات (page setup) ، یعنی حواشی بالا،‌ پایین، چپ و راست (پیش‌فرض طبق اعداد تزهای دانشگاه شریف).
% چهارم: تغییر فرمت header ، طوری که یک خط در آن کشیده شود و تیتر فصل و شماره صفحه بالای خط نوشته شود (fancy header).
% پنجم: نحوه تیتر گذاری شروع هر فصل تغییر داده شد: بالای آن با خط توخالی بنویسد (مثلا) «فصل ۱» و بعد زیر آن تیتر فصل را بنویسد و همه هم در وسط باشد. به نظرم از پیش‌فرض  report (که در واقع برای انگلیسی است) زیباتر است، لااقل برای تز.
% ششم: تغییر نحوه شماره‌گذاری فصلها و فرمولها و غیره، طوری که مثلا شماره فرمول به صورت (۱-۵) نوشته شود نه (۱.۵).
% هفتم: تغییر اینکه بجای «لیست اشکال» و «لیست جداول» بنویسد «فهرست اشکال» و «فهرست جداول».
% هشتم: تعریف دو محیط note (توجه) و experiment (آزمایش).
% نهم: تعریف محیطهایی مثل قضیه، تعریف و غیره (که نهایتا کامنت شده‌اند و غیرفعال هستند و کاربر اگر بخواهد می‌تواند آنها را فعال کند).
%
% نویسنده: مسعود بابایی‌زاده

% نسخه 1.0.9
% تاریخ: ۵ اسفند ۱۴۰۰
% تغییرات:
%   زی‌پرشینی که در دی وی دی texlive2021 بود (زیرپرشین نسخه 23.1) باگهای متعددی پیدا کرده است که تغییراتی که در این نسخه فایل حاضر داده شده است، همگی برای این است که فایل بتواند با این نسخه زی‌پرشین کار کند. اگر در نسخه‌های بعدی زی‌پرشین این مشکلات رفع شده باشد، بهتر است به نسخه قبلی (نسخه 1.0.8) برگشت. مشکلات زی‌پرشین جدید و تغییرات فایل حاضر برای رفع آنها عبارتند از:
%
%      ۱) این زی‌پرشین با فونت XB Titre Shadow که من برای نوشتن کلمه فصل (مثلا نوشتن کلمه «فصل ۱» بالای عنوان فصل اول) استفاده کرده بودم درست کار نمی‌کرد و آن را به صورت منفصل می‌نوشت («ف‌ص‌ل ۱»). پس این فونت را با IRSina جایگزین کردم (البته فونتهای خیلی زیادی امتحان کردم، بخصوص توخالی. ولی هیچدام جالب نبودند. شاید Mj_Nazila 3D Shadow بد نبود، ولی اعدادش بجای فارسی، عربی بودند). به هر حال به نظرم همان XB Titre Shadow خیلی بهتر است و چون ظاهرا این مشکل از خود زی‌پرشین است و نه فونت، در نسخه‌های آینده زی‌پرشین دقت شود که اگر این مشکل درست شده بود، مجددا به فونت XB Titre Shadow برگردیم.
%
%      ۲) صدا زدن دستور \setmathdigitfont در زی‌پرشین جدید منجر به یک خطای کمپایل می‌شود! راه حل پیشنهاد شده در لینک زیر در فایل حاضر استفاده شد:
%    http://qa.parsilatex.com/35710/__xepersian_mathsdigitspec_primitive_font_char_if_exist
%
%      ۳) محیط جدید document که با دستور زیر تعریف کرده بودم، در زی‌پرشین جدید این خطای کمپایل را می‌داد که \begin{olddocument} یا \end{document} تمام شده است!
%    \renewenvironment{document}{\begin{olddocument}\setlength{\baselineskip}{1.53\baselineskip}}{\end{olddocument}}
% که به وضوح باگ زی‌پرشین جدید است. در نتیجه دستور بالا را جایگرین کردم با دستور زیر:
%    \renewenvironment{document}{\olddocument\setlength{\baselineskip}{1.53\baselineskip}}{\endolddocument}


%%%%%%%%%%%%%%%%%%%%%%%%%%%%%%%%%%%%%%%%%%
% نسخه 1.0.8
% تاریخ: ۵ اسفند ۱۴۰۰
% تغییرات:
% تغییرات:
%           ۱) فونتهای پیش‌فرض متن که قبلا IRLotus بودند با IRLotusICEE (که نسخه اصلاح‌شده‌ای از IRLotus است که خودم برای کنفرانس برق ایران سال ۲۰۱۵ درست کرده بودم)، تغییر داده شد (ابتدا MBZLotus را هم اضافه کرده بودم، ولی در محیط itemize کار نمی‌کند و کاراکتر بولت آیتم را ندارد. در نتیجه حذف شد).
%
%           ۲) دستور \setdigitfont با \setmathdigitfont جایگزین شد تا با زی‌پرشین جدید درست کار کند و در عنوان فصلها، مثلا در  عبارت «فصل ۱»، عدد ۱ هم با همان فونت کلمه «فصل» (یعنی توخالی) نوشته شود، نه توپر.
%		   
%		   ۳) شماره‌گذاری فرمولها و زیرفصلها و غیره درست شد تا در زی‌پرشین جدید درست کار کند (وگرنه مثلا شماره فرمول ۲ فصل ۱، بجای (۲-۱) نوشته می‌شد (۱-۲)).

%%%%%%%%%%%%%%%%%%%%%%%%%%%%%%%%%%%%%%%%%%

% نسخه 1.0.7
% تاریخ: ۲۲ مهر ۱۳۹۳
% تغییرات: ۱) برای کارکردن با texlive2014 و زی‌پرشین جدید:
%               - فونتها با IRLotus جایگزین شدند => فقط با texlive2013 و texlive2014 کار می‌کند (با 2012
%                  کار نمی‌کند، چون کشیده و کاراکتر نیمفاصله را در خروجی بین دو عدد "" می‌گذارد!).
%               - بجای داشتن یک فایل MBZ_ThesisSettings.sty، دو فایل MBZ_ThesisSettings1.tex و
%                  MBZ_ThesisSettings2.tex داریم، که اولی شامل تنظیماتی است که قبل از لود کردن زی‌پرشین
%                  باید انجام شوند و دومی تنظیماتی که بعد از آن باید انجام شوند.
%          ۲) تعریف فونتهایی که در فایل Jeld.tex بود به همان MBZ_ThesisSettings2.tex منتقل شدند. اکنون همه
%             فونتها در همین فایل تعریف می‌شوند.
%          ۳) شماره گذاری حرفی طوری تغییر یافت که بجای «آ-ب-ج-د...» داشته باشیم «الف-ب-ج-د...» (با استفاده
%             از دستورات دکتر امین‌طوسی در تمپلیت دانشگاه خودشان.


%%%%%%%%%%%%%%%%%%%%%%%%%%%%%%%%%%%%%%%%%%
% نسخه‌های قبلی:
%%%%%%
% نسخه 1.0.6
% تاریخ: ۲۵ بهمن ۱۳۸۸
% تغییرات: ۱) تعریف محیط‌های definition، theorem، و ...
%           ۲) نحوه شماره‌گذاری table را فراموش کرده بودم که مثل figure درست کنم، اصلاح شد.
% ------>   3) bbding و محیط experiment درست شوند.  <-------
%           ۴) \headheigh هم تنظیم شد (برابر ۱۳ پوینت). قبلا اصلا دست نزده بودم، ظاهر پیش‌فرض آن ۱۲ بوده و fancyhdr وارنینگ می‌داد.
%%%%%%
% نسخه 1.0.5
% تاریخ: ۲۰ بهمن ۱۳۸۸
% تغییرات: در فهرست عنوان فصلها را بجای اینکه \Large بنویسم، کمی کوچکتر یعنی \large نوشتم (و عنوان \part ها کماکان \Large است).
%%%%%%
% نسخه 1.0.4
% تاریخ: ۱۸ بهمن ۱۳۸۸
% تغییرات: فونت تیتر فصلها را عوض کردم و طوری نوشتم که دو فونت متفاوت بگیرد، یکی برای شماره فصل (مثلا «فصل اول») و دیگری برای عنوان فصل (مثلا «مقدمه»). همچنین اولی را} XB Titre Shadow و دومی را Titre گذاشتم.

%%%%%%
% نسخه 1.0.3
% تاریخ: ۱۶ بهمن ۱۳۸۸
% تغییرات: نحوه نوشتن part را هم در فهرست عوض کردم، تا در فارسی زیبا شود. همچنین در فهرست اندازه تیتر فصلها را نیز از \Large به \large کاهش دادم، و اندازه تیتر part را به \Large نوشتم.
%%%%%%
% نسخه 1.0.2
% تاریخ: ۱۵ بهمن ۱۳۸۸
% تغییرات: فاصله فرمولها از متن بیشتر شد (به نظرم این فاصله را زیادی کم کرده بودم). کامنتهای بالا هم ترتیبشان تغییر کرد و تا حدودی به ترتیب اولویت شدند.

%%%%%%
% نسخه: ۱.۰.۱
% تاریخ: ۲۵ دیماه ۱۳۸۸ برابر با ۱۵ ژانوبه ۲۰۱۰
% تغییرات: تغییر داده شده نسبت به نسخه 1.0.0 فقط زیادکردن \parindent است. مقدار پیش‌فرض قبلی، برای متون انگلیسی است و در متون فارسی زیادی کوچک است.

%%%%%%%%%%%%%%%%%%%%%%%%%%%%%%%%%%%%%%%%%%
% Font settings

%\defpersianfont\chaptertitlenumberfont[Scale=1.31]{B Koodak Outline}
%\defpersianfont\chaptertitlenumberfont[Scale=1]{XB Titre Shadow}
\defpersianfont\chaptertitlenumberfont[Scale=1]{IRSina.ttf}
\defpersianfont\chaptertitlefont[Scale=1]{IRTitr.ttf}

% فونت اصلی متن فونت «یاس» است که توسط کاربران xepersian از اصلاح XB Yas به دست آمده است.
%\settextfont[Scale=1.1]{IRLotus}
%\setlatintextfont[Scale=1]{Times New Roman}
%\setonlydigitfont[Scale=1.1]{Yas} %This is not a xepersian command. It is a command provided by setdigitfont.sty, written by Massoud Babaie-Zadeh
%\setdigitfont[Scale=1.1]{IRLotus}

% استفاده از حالت خوابیده به چپ قلم «یاس» مشابه با حالت ایرانیک در فارسی‌تک
%\setiranicfont[Scale=1.1]{IRLotus-Italic}

% 12 point "IRLotus" is small. Use 13.1 point (which gives a size similar to FarsiTeX)
\settextfont[ BoldFont={IRLotusICEE_Bold.ttf}, BoldItalicFont={IRLotusICEE_BoldIranic.ttf}, ItalicFont={IRLotusICEE_Iranic.ttf},Scale=1.31]{IRLotusICEE.ttf}%{IRZar.ttf}
%\setdigitfont[Scale=1.31]{IRLotusICEE} % replaced by the next command since texlive2018:
% دستور بالا از نسخه‌ای به بعد زی‌پرشین (لااقل از texlive2018 به بعد) باعث میشود که همه عددها با فونت بالا جایگزین شود که مطلوب نیست. مثلا در عنوان فصل که نوشته شده «فصل ۱» هم کلمه «۱» با این فونت (و در نتیجه توپر) نوشته میشود در حالیکه «فصل» با فونت توخالی بوده. و برای اینکه فقط اعداد داخل فرمولها فارسی شوند، دستور بالا با دستور جدید \setmathdigitfont در زیر جایگزین شد:

% شروع دستوراتی برای رفع یک باگ در زی‌پرشین 23.1 که در دی وی دی texlive2021 است. منبع: http://qa.parsilatex.com/35710/__xepersian_mathsdigitspec_primitive_font_char_if_exist
\ExplSyntaxOn
\cs_set_eq:NN
\etex_iffontchar:D
\tex_iffontchar:D
\cs_undefine:N \c_one
\int_const:Nn \c_one { 1 }
\ExplSyntaxOff
% پایان دستورات رفع باگ بالا در زی‌پرشین 23.1
\setmathdigitfont[Scale=1.31]{IRLotusICEE.ttf}
%\setlatintextfont[Scale=1]{Times New Roman}
\setiranicfont[Scale=1.31]{IRLotusICEE_Iranic.ttf}				% ایرانیک، خوابیده به چپ

% Fonts on JELD (title page)
\defpersianfont\mainjeldfont[ BoldFont={IRLotusICEE_Bold.ttf}, BoldItalicFont={IRLotusICEE_BoldIranic.ttf}, ItalicFont={IRLotusICEE_Iranic.ttf},Scale=1.31]{IRLotusICEE.ttf}
\defpersianfont\fieldnamefont[Scale=1.2]{BKoodakO.ttf}
\defpersianfont\ITRCthanksfont[Scale=1]{BTraffic.ttf}


%%%%%%%%%%%%%%%%%%%%%%%%%%%%%%%%%%%%%%%%%%
% دستور بعدی برای زیادکردن فاصلهٔ بین خطوط نوشته شده بود:
%\renewcommand{\baselinestretch}{2.3}
%\setstretch{2.3}
%\linespread{2.3}

% در اینصورت مشکل آن است که نه تنها فاصله بین خطوط متون فارسی را زیاد کنیم، فاصله خطوط بین فرمولها (مثلا در ماتریسها) هم زیاد می‌شود (که عدد ۲.۳ بالا زیادی بزرگ است). بعد برای حل این مشکل ابتدا سعی کردم محیط equation را طوری تغییر دهیم که اول فاصله خطوط را کمتر کند، بعد equation قدیم را صدا بزند. من متناسب با فاصله خطوط فارسی برابر 2.3 برابر (عدد بالا) فاصله فرمولها برابر 1.5 به نظرم خوب آمد، پس اول همه این دستورات را امتحان کردم:
%\let\oldequation=\equation
%\let\endoldequation=\endequation
%\renewenvironment{equation}{\vspace{-1em}\begin{spacing}{1.5}\begin{oldequation}}{\end{oldequation}\end{spacing}\vspace{-1em}}

%\renewenvironment{equation}{\begingroup\setstretch{1}\begin{oldequation}}{\end{oldequation}\setstretch{2.3}\endgroup}

%\renewenvironment{equation}{\vskip -\parskip\vskip -\baselineskip\begin{spacing}{1.5}\begin{oldequation}}{\end{oldequation}\end{spacing}\vskip -\parskip\vskip -\baselineskip}

%\renewenvironment{equation}{\begin{spacing}{1.5}\begin{oldequation}}{\end{oldequation}\end{spacing}}

%\let\old\baselinestretch=\baselinestretch\renewcommand{\baselinestretch}{2.3}
%\renewenvironment{equation}{\begingroup\def\baselinestretch{1.5}\begin{oldequation}}{\end{oldequation}\def\baselinestretch{2.3}\endgroup}

 % اما هر کدام مشکلی داشت. بهتر از همه همان اولی بود که مشکل آن این بود که پاراگراف بعد از فرمولها همیشه indent می‌شود.

% در نهایت راه حل بهتری پیدا کردم. متوجه شدم که زیاد کردن \baselineskip بر خلاف \baselinestreatch روی محیط ریاضی تاثیری ندارد. اما \baselineskip را باید بعد از \begin{document} زیاد کرد. با توجه به اینکه singlespace برای فرمولهای ریاضی در متن فارسی ما زیادی کوچک است،‌ پس برای آنکه طبق اعداد بالا فاصله خطوط در فرمولهای 1.5 برابر و در متن فارسی 2.3 برابر باشد، لازم است که طبق دستور زیر \baselinestreatch برابر 1.5 قرار داده شود و سپس درون متن و بعد از  \begin{document} باید \baselineskip را 2.3/1.5=1.53 برابر نمود. یعنی:

\renewcommand{\baselinestretch}{1.5}
%\setlength{\baselineskip}{1.53\baselineskip}   ->  This is inside the text and right after \begin{document}
%برای آنکه کاربر مجبور نباشد دستور بالا را دستی بعد از begin document اضافه کند، دستورات زیر را می‌نویسیم:
\let\olddocument=\document
\let\endolddocument=\enddocument
%\renewenvironment{document}{\begin{olddocument}\setlength{\baselineskip}{1.53\baselineskip}}{\end{olddocument}}
\renewenvironment{document}{\olddocument\setlength{\baselineskip}{1.53\baselineskip}}{\endolddocument}
%در اینصورت متوجه شدم که فاصله فرمولها با متن کمی زیاد می‌شود که آن را نیز با دستورات زیر می‌توان حل کرد (عدد 0.5em در نسخه 1.0.2 گذاشته شد، قبلا بالا را -0.6em و پایین را 0.7em گذاشته بودم که به نظرم فرمولها زیادی به متن نزدیک می‌شد):
\let\oldequation=\equation
\let\endoldequation=\endequation
\renewenvironment{equation}{\vspace{-0.2em}\begin{oldequation}}{\vspace{-0.2em}\end{oldequation}\ignorespacesafterend}


% همچنین متوجه شدم که با اعداد بالا در فهرست مطالب و فهرست اشکال و جداول نیز فاصله خطوط زیاد است. که به صورت زیر می‌توان اصلاح کرد (یعنی برای آنها baselineskip را مجددا به عدد قبلی برگرداند، یعنی در معکوس 1.53 که برابر 0.65 می‌شود ضرب کرد):
\let\oldtableofcontents=\tableofcontents
\renewcommand{\tableofcontents}{\begingroup\setlength{\baselineskip}{0.65\baselineskip}\oldtableofcontents\endgroup}

\let\oldlistoffigures=\listoffigures
\renewcommand{\listoffigures}{\begingroup\setlength{\baselineskip}{0.65\baselineskip}\oldlistoffigures\endgroup}

\let\oldlistoftables=\listoftables
\renewcommand{\listoftables}{\begingroup\setlength{\baselineskip}{0.65\baselineskip}\oldlistoftables\endgroup}

% همچنین متوجه شدم که دستور با اعداد بالا، فاصله خطوط در یک متن انگلیسی زیادی  (مثلا فهرست مراجع) بزرگ است. برای متن انگلیسی حتی مشابه بالا ضرب کردن baselineskip در 0.65 هم به نظرم کافی نبود و در پایین با تغییر تعریف latin آن را در 0.55 ضرب کرده‌ام:
\let\oldlatin=\latin
\let\endoldlatin=\endlatin
\renewenvironment{latin}{\begin{oldlatin}\setlength{\baselineskip}{0.55\baselineskip}}{\end{oldlatin}}

%%%%%%%%%%%%%%%%%%%%%%%%%%%%%%%%%%%%%%%%%%
% اصلاحات مربوط به فهرست مطالب (تغییرات الف تا د زیر):
% الف) یک تغییر، کم‌کردن فاصله خطوط است که در بالا با اصلاح تعریف \tableofcontents انجام شده بود.
% ب) دستور زیر تغییر داده شده است: تغییر اول برای آن است که قبل از نوشتن تیتر فصلها کمی فاصله بیشتر نسبت به فاصله خطوط در زیرفصلها بگذارد و دستور دوم برای آن است که در فهرست مطالب، اسم فصلها را درشت‌تر بنویسد.

\makeatletter

\renewcommand{\l@chapter}[2]{%
  \ifnum \c@tocdepth >\m@ne
    \addpenalty{-\@highpenalty}%
    \vskip 2.0em \@plus\p@ % This line is modified from the default definition (in "report.cls" file): \vskip 1.0em was changed to \vskip 2.0em (Done by Massoud Babaie-Zadeh, 14 Jan 2010)
    \setlength\@tempdima{1.5em}%
    \begingroup
      \parindent \z@ \rightskip \@pnumwidth
      \parfillskip -\@pnumwidth
      \leavevmode \bfseries
      \advance\leftskip\@tempdima
      \hskip -\leftskip
      \begingroup\large #1\endgroup\nobreak\hfil \nobreak\hb@xt@\@pnumwidth{\hss #2}\par % This line is modified from the default definition: #1 has been replaced by \begingroup\Large #1\endgroup (Done by Massoud Babaie-Zadeh, 14 Jan 2010)
      \penalty\@highpenalty
    \endgroup
  \fi}

\makeatother
  
% ج) دستورات زیر نحوه نوشتن part را در فهرست عوض می‌کند. سه دستور تغییر داده شده است: اول \partname است که پیش‌فرض xepersian تغییر داده شده تا بجای «بخش» بنویسد «قسمت». دو دستور دیگر از خود لاتک هستند (فایل report.cls) که نحوه نوشتن را عوض می‌کنند.
  
\def\partname{قسمت}  

\makeatletter
% Many parts of the following command are changed.
\renewcommand*\l@part[2]{%
  \ifnum \c@tocdepth >-2\relax
    \addpenalty{-\@highpenalty}%
    \addvspace{2.25em \@plus\p@}%
    \setlength\@tempdima{1.5em}%
    \begingroup
      \parindent \z@ \rightskip \@pnumwidth
      \parfillskip -\@pnumwidth
      \leavevmode
      \advance\leftskip\@tempdima
      \hskip -\leftskip
	   \Large \bfseries #1\hfil \hb@xt@\@pnumwidth{\hss #2}\par
       \nobreak
         \global\@nobreaktrue
         \everypar{\global\@nobreakfalse\everypar{}}%
    \endgroup
  \fi}
  
\def\@part[#1]#2{%
    \ifnum \c@secnumdepth >-2\relax
      \refstepcounter{part}%
	  \addcontentsline{toc}{part}{\partname\ \thepart\unskip: #1}%  -> This is the only part modified by Massoud Babaie-Zadeh (on 5 Feb 2010). The original was:
      %\addcontentsline{toc}{part}{\thepart\hspace{1em}#1}%
    \else
      \addcontentsline{toc}{part}{#1}%
    \fi
    \markboth{}{}%
    {\centering
     \interlinepenalty \@M
     \normalfont
     \ifnum \c@secnumdepth >-2\relax
       \huge\bfseries \partname\nobreakspace\thepart
       \par
       \vskip 20\p@
     \fi
     \Huge \bfseries #2\par}%
    \@endpart}
\def\@spart#1{%
    {\centering
     \interlinepenalty \@M
     \normalfont
     \Huge \bfseries #1\par}%
    \@endpart}
\def\@endpart{\vfil\newpage
              \if@twoside
               \if@openright
                \null
                \thispagestyle{empty}%
                \newpage
               \fi
              \fi
              \if@tempswa
                \twocolumn
              \fi}
  
  
%  د) دستور زیر برای آن است که در فهرست مطالب بعد از هر section یک سطر خالی بگذارد تا قسمتهای مختلف کمی از هم مجزا بشوند.
\let\oldl@section=\l@section
\renewcommand{\l@section}{\vspace{1em}\oldl@section}

\makeatother



% خطهای زیر تنها برای کشیده‌شدن یک خط در header اضافه شده‌اند. در
% نتیجه، پکیج ffancyhdr نیز باید اضافه شود.
\lhead{\thepage}
\chead{}
\rhead{\leftmark}
\lfoot{}
\cfoot{}
\rfoot{}
\pagestyle{fancy}
\renewcommand{\chaptermark}[1]{\markboth{\chaptername\ \thechapter:\ #1}{}}

% End of Page Setup
%%%%%%%%%%%%%%%%%%%%%%%%%%%%%%%%%%%%%%%%%%
% دستور زیر برای زیادکردن تورفتگی اول هر پاراگراف است. مقدار پیش‌فرض قبلی، برای متون انگلیسی است و برای متون فارسی زیادی کوچک است.

\parindent=1cm


%%%%%%%%%%%%%%%%%%%%%%%%%%%%%%%%%%%%%%%%%%
% دستور زیر برای تنظیم تیتر شروع هر فصل نوشته شده است. دستور اول فونت نوشتن عبارت «فصل ؟» را تعیین می‌کند و دستور بعدی، تنظیم خود تیتر است. (این دستور در واقع در فایل report.cls در LateX2e و در rep10.sty, rep11.sty, rep12.sty در LaTeX2.09 تعریف شده است):

\makeatletter
 
\def\@makechapterhead#1{%
  \vspace*{50\p@}%
  {\parindent \z@ \centering \normalfont
    \ifnum \c@secnumdepth >\m@ne
        {\Large\chaptertitlenumberfont  \@chapapp\space \thechapter}
        \par\nobreak
        \vskip 20\p@
    \fi
    \interlinepenalty\@M
    \Huge\chaptertitlefont\bfseries #1\par\nobreak
    \vskip 60\p@
  }}

% و برای اینکه همین اتفاق در مورد دستور chapter* هم بیفتد (که مثلا در نوشتن تیتر «فهرست مطالب» و «لیست اشکال» و «لیست جداول» و غیره استفاده می‌شود)
\def\@makeschapterhead#1{%
  \vspace*{50\p@}%
  {\parindent \z@ \centering
    \normalfont
    \interlinepenalty\@M
    \Huge \chaptertitlefont\bfseries  #1\par\nobreak
    \vskip 60\p@
  }}
\makeatother

%%%%%%%%%%%%%%%%%%%%%%%%%%%%%%%%%%%%%%%%%%
% فرمان بعدی بدین منظور نوشته شده است که شماره فرمولها بصورت
% (شماره فصل-شماره فرمول) نوشته شود.
\renewcommand{\theequation}{\thechapter-\arabic{equation}}
% به همین ترتیب در مورد سایر شماره‌ها (شکلها، بخش‌ها، ...):
\renewcommand{\thefigure}{\thechapter-\arabic{figure}}
\renewcommand{\thetable}{\thechapter-\arabic{table}}
\renewcommand{\thesection}{\thechapter-\arabic{section}}
\renewcommand{\thesubsection}{\thesection-\arabic{subsection}}
\renewcommand{\thesubsubsection}{\thesubsection-\arabic{subsubsection}}
\renewcommand{\theparagraph}{\thesubsubsection-\arabic{paragraph}}
\renewcommand{\thesubparagraph}{\theparagraph-\arabic{subparagraph}}

%%%%%%%%%%%%%%%%%%%%%%%%%%%%%%%%%%%%%%%%%%
% دو دستور بعدی برای این نوشته شده‌اند که به جای «لیست اشکال» و «لیست جداول» که مقدار
% پیش‌فرض فارسی‌تک و xepersian است، نوشته شود  «فهرست اشکال» و «فهرست جداول». 
% همچنین دستور \hfil برای آن است که وسط نوشته شود.
\def\listfigurename{فهرست اشکال}
\def\listtablename{فهرست جداول}


%%%%%%%%%%%%%%%%%%%%%%%%%%%%%%%%%%%%%%%%%%
% این قسمت برای فراهم آوردن کارکتر «پیچ خطرناک» و تهیه محیط note نوشته
% شده است:
\font\manfnt=manfnt
\newlength{\dbendheight}
\def\lhdbend{{\manfnt \char126}}
\newcommand{\textlhdbend}{\raisebox{10pt}{\lhdbend}}
\newenvironment{note}{\vspace{1.5em}\noindent\textlhdbend {\bf توجه:}\trafficfont%
}{\vspace{1.5em}}

%%%%%%%%%%%%%%%%%%%%%%%%%%%%%%%%%%%%%%%%%%
% برای آنکه در شماره‌گذاری حرفی و ابجد به جای آ از الف استفاده شود (این دستورات از تمپلیت تهیه شده توسط دکتر امین‌طوسی برا پایان‌نامه‌های دانشگاه حکیم سبزواری برداشته شده است):

\makeatletter

 \def\abj@num@i#1{%
   \ifcase#1\or الف\or ب\or ج\or د%
            \or ه‍\or و\or ز\or ح\or ط\fi
   \ifnum#1=\z@\abjad@zero\fi}   
  
   \def\@harfi#1{\ifcase#1\or الف\or ب\or پ\or ت\or ث\or
 ج\or چ\or ح\or خ\or د\or ذ\or ر\or ز\or ژ\or س\or ش\or ص\or ض\or ط\or ظ\or ع\or غ\or
 ف\or ق\or ک\or گ\or ل\or م\or ن\or و\or ه\or ی\else\@ctrerr\fi}
 
\makeatother


%%%%%%%%%%%%%%%%%%%%%%%%%%%%%%%%%%%%%%%%%%
% این قسمت برای فراهم کردن کاراکتر «مداد کوچک» و تهیه محیط experiment
% نوشته شده است.
%\font\bbding=bbding at 15pt
%\def\smallpencil{{\bbding \char28}}
%\def\expfont{\traffic}
%\newenvironment{experiment}{\vspace{1em}\smallpencil {\bf آزمایش:}%
%\expfont}%
%{\rm  \vspace{1em}}

%%%%%%%%%%%%%%%%%%%%%%%%%%%%%%%%%%%%%%%%%%
% این قسمت تعریف محیطهای lemma، theorem و غیره است.

\newtheorem{definition}{تعریف}[chapter]
\renewcommand{\thedefinition}{{\arabic{definition}-\thechapter}}

\newtheorem{property}{خاصیت}[chapter]
\renewcommand{\theproperty}{{\arabic{property}-\thechapter}}

\newtheorem{theorem}{قضیه}[chapter]
\renewcommand{\thetheorem}{{\arabic{theorem}-\thechapter}}

\newtheorem{lemma}{لم}[chapter]
\renewcommand{\thelemma}{{\arabic{lemma}-\thechapter}}

\newtheorem{corollary}{نتیجهٔ فرعی}
\renewcommand{\thecorollary}{{\arabic{corollary}-\thechapter}}

%\newtheorem{predefinition}{تعریف}[chapter]
%\renewcommand{\thepredefinition}{{\arabic{predefinition}-\thechapter}}
%\newenvironment{definition}{\begin{predefinition}{\hspace{-0.4em}{\bf.}}}
%                          {\end{predefinition}}

%\newtheorem{preproperty}{خاصیت}[chapter]
%\renewcommand{\thepreproperty}{{\arabic{preproperty}-\thechapter}}
%\newenvironment{property}{\begin{preproperty}{\hspace{-0.4em}{\bf.}}}
%                          {\end{preproperty}}

%\newtheorem{pretheo}{قضیه}[chapter]
%\renewcommand{\thepretheo}{{\arabic{pretheo}-\thechapter}}
%\newenvironment{theorem}{\begin{pretheo}{\hspace{-0.4em}{\bf.}}}
%                          {\end{pretheo}}

%\newtheorem{prelemma}{لم}[chapter]
%\renewcommand{\theprelemma}{{\arabic{prelemma}-\thechapter}}
%\newenvironment{lemma}{\begin{prelemma}{\hspace{-0.4em}{\bf.}}}
%                          {\end{prelemma}}

%\newtheorem{precorollary}{نتیجهٔ فرعی}
%\renewcommand{\theprecorollary}{{\arabic{precorollary}-\thechapter}}
%\newenvironment{corollary}{\begin{precorollary}{\hspace{-0.4em}{\bf.}}}
%                          {\end{precorollary}}


%\newenvironment{proof}{\vspace{0.5em} {\bf اثبات.} 
%\rm  }{\hfill{$\Box$} \vspace{2em}}

%\newenvironment{example}{\vspace{0.5em}{\bf مثال.} 
%\small  }{\vspace{1em}}

%%%%%%%%%%%%%%%%%%%%%%%%%%%%%%%%%%%%%%%%%%
 % Settings after loading xepersian (written by Massoud Babaie-Zadeh)
% Bold Symbols
\newcommand{\yb}{{\mathbf y}}
\newcommand{\sbb}{{\mathbf s}}
\newcommand{\xb}{{\mathbf x}}
\newcommand{\eb}{{\mathbf e}}
\newcommand{\tb}{{\mathbf t}}
\newcommand{\wb}{{\mathbf w}}
\newcommand{\kb}{{\mathbf k}}
\newcommand{\Tb}{{\mathbf T}}
\newcommand{\Ab}{{\mathbf A}}
\newcommand{\Eb}{{\mathbf E}}
\newcommand{\Bb}{{\mathbf B}}
\newcommand{\Mb}{{\mathbf M}}
\newcommand{\Nb}{{\mathbf N}}
\newcommand{\Bbh}{{\mathbf{\hat B}}}
\newcommand{\ybh}{{\mathbf{\hat y}}}
\newcommand{\xbh}{{\mathbf{\hat x}}}
\newcommand{\Rbh}{{\mathbf{\hat x}}}
\newcommand{\Bbs}{{\mathbf{\scriptsize B}}}
\newcommand{\Db}{{\mathbf D}}
\newcommand{\Ib}{{\bf I}}
\newcommand{\Rb}{{\mathbf R}}
\newcommand{\Ei}{{\mbox{$\cal E$}}}
\newcommand{\RR}{{\Bbb R}}
\newcommand{\bijk}{b_{ij}^{(k)}}
\newcommand{\tetab}{{\mbox{\boldmath $\theta$}}}
\newcommand{\phib}{{\mbox{\boldmath $\phi$}}}
\newcommand{\psib}{{\mbox{\boldmath $\psi$}}}
\newcommand{\mub}{{\mbox{\boldmath $\mu$}}}
\newcommand{\betab}{{\mbox{\boldmath $\beta$}}}
\newcommand{\Deltab}{{\mbox{\boldmath $\Delta$}}}
\newcommand{\deltab}{{\mbox{\boldmath $\delta$}}}
\newcommand{\etab}{{\mbox{\boldmath $\eta$}}}
\newcommand{\alphab}{{\mbox{\boldmath $\alpha$}}}
\newcommand{\gammab}{{\mbox{\boldmath $\gamma$}}}
\newcommand{\xib}{{\mbox{\boldmath $\xi$}}}
\newcommand{\Ub}{{\mathbf U}}
\newcommand{\Vb}{{\mathbf V}}
\newcommand{\Xb}{{\mathbf X}}
\newcommand{\Sb}{{\mathbf S}}

% Calligraphic symbols
\newcommand{\cC}{{\cal C}}
\newcommand{\cL}{{\cal L}}
\newcommand{\cU}{{\cal U}}
\newcommand{\cE}{{\cal E}}
\newcommand{\cG}{{\cal G}}

% News Commands
\newcommand{\eqnarrayeq}{\!\!\!=\!\!\!}
\newcommand{\lk}{ \left\{ }
\newcommand{\rk}{ \right\} }
\newcommand{\la}{ \left\langle }
\newcommand{\ra}{ \right\rangle }
\newcommand{\tr}{ {\mbox{trace}} }
\newcommand{\rond}[1]{\frac{\partial}{\partial #1}}
\newcommand{\varrond}[2]{\frac{\partial #1}{\partial #2}}
\newcommand{\mb}{\left(}   % Matrix Begin
\newcommand{\me}{\right)}  % Matrix End
\newcommand{\snr}{\mbox{SNR}}
\newcommand{\argmin}{\mathop{\engmbox{\rm argmin}}}
\newcommand{\lmat}{\left[} % Left matrix
\newcommand{\rmat}{\right]}% Right matrix
\newcommand{\myrightarrow}{\mathop{\longrightarrow}}
\newcommand{\spl}{\mbox{sp}}
 % برای دستورات جدیدی که خودتان می‌خواهید تعریف کنید.


%%%%%%%%%%%%%%%%%%%%%%%%%%%%%%%%%%%%%%%%%%%%%%%%%%%%
\begin{document}

% صفحات اولیه تز (جلد، بسم‌ا...، توضیح (کامنت)، تصویب‌نامه، تشکر، چکیده) %%%%%%%%%%%%%%%
\include{jeld}
\thispagestyle{empty}
\vspace*{5cm}

\begin{latin}
\centering
\includegraphics{Besmellah.eps}
\end{latin}

\thispagestyle{empty}
\phantom{a}
\vfill
\begin{center}
\fbox{\begin{minipage}{9cm}
\begin{center}
\begin{minipage}{8cm}
\vspace*{7mm}

\mainjeldfont

{\begin{center}\large\textbf{توجه}\end{center}}

این پروژه بر اساس قرارداد شماره (............) از حمایت مالی مرکز تحقیقات مخابرات ایران 
برخوردار شده است.

\vspace*{7mm}
\end{minipage}
\end{center}
\end{minipage}}
\end{center}
\vspace{5cm}\vfill
%\begin{center}


\thispagestyle{empty}
\begin{center}
\vskip -2cm

{\Large\textbf{بسمه تعالی}}
\vskip 1cm
{\Large\textbf{دانشگاه صنعتی شریف} \\ \textbf{دانشکده مهندسی برق}}
\vskip 1.5cm
{\Large\textbf{پایان‌نامه کارشناسی ارشد}} \par 
\end{center}
\vskip 2cm

{عنوان: \textbf{طراحی و پیاده‌سازی یک پلتفرم MLOps به‌صورت ابری مبتنی‌بر GPU}}

{نگارش: \textbf{ابوالفضل یاریان}}


\vskip 2cm
{\Large\textbf{اعضا هیات  داوران:}} \par
\vskip 5mm
{\textbf{ دکتر متین هاشمی  } \hfill 
{ امضاء:........................... } \hspace{1cm} \vspace{2mm}\par
{\textbf{ دکتر \dots   } \hfill 
{ امضاء:........................... } \hspace{1cm} \vspace{2mm}\par
 {\textbf{ دکتر \dots   } \hfill 
{ امضاء:........................... } \hspace{1cm} \vspace{2mm}\par
{\textbf{ دکتر \dots   } \hfill 
{ امضاء:........................... } \hspace{1cm} \vspace{2mm}\par
{\textbf{ دکتر \dots   } \hfill 
{ امضاء:........................... } \hspace{1cm} \vspace{2mm}\par

\vskip 1.5cm
{ تاریخ: } \par 


\thispagestyle{empty}

\vspace*{1cm}

\centerline{\Large\textbf{تقدیم و قدردانی}}
\vskip 1.8cm

{
این پایان‌نامه را به پدر و مادر عزیزم تقدیم می‌کنم که همواره حامی و مشوق من بوده‌اند. از استاد راهنمای گرامی، دکتر متین هاشمی، بابت راهنمایی‌ها و حمایت‌های بی‌دریغشان صمیمانه تشکر می‌کنم. همچنین، قدردانی ویژه‌ای دارم از دکتر طاها قاسمی به خاطر مشاوره‌ها و راهنمایی‌های ارزشمندشان در طول این مسیر و از تمامی دوستان و همکارانی که با همدلی و همراهی مرا در این مسیر یاری دادند نیز بی‌نهایت سپاسگزارم.
}

%\begin{itemize}
% \item
% \item 
%\end{itemize}




\thispagestyle{empty}
\phantom{a}
\vfil
%\voffset=1cm

\begin{center}
\begin{minipage}{0.9\textwidth}


\noindent \textbf{چکیده:} 
\vskip 2mm \par

در عصر حاضر، هوش مصنوعی و یادگیری ماشین به یکی از حیاتی‌ترین و پرکاربردترین فناوری‌ها در صنایع مختلف تبدیل شده‌اند. این فناوری‌ها به شرکت‌ها و سازمان‌ها این امکان را می‌دهند تا بهینه‌سازی فرآیندها، پیش‌بینی روندها، و کشف الگوهای پنهان در داده‌ها را با دقت و سرعت بالا انجام دهند. با این حال، بهره‌گیری کامل از قابلیت‌های هوش مصنوعی و یادگیری ماشین نیازمند پیاده‌سازی مؤثر و کارآمد مدل‌های یادگیری ماشین در محیط‌های تولیدی است. \lr{MLOps}، که ترکیبی از مفاهیم \lr{DevOps} و \lr{ML} است، به مدیریت چرخه عمر مدل‌های یادگیری ماشین از توسعه تا استقرار و نگهداری کمک می‌کند. در این تحقیق، به دلیل تحریم‌های بین‌المللی و محدودیت‌های دسترسی به سرویس‌های خارجی مانند \lr{Google Vertex AI}، \lr{Amazon SageMaker} و \lr{Microsoft Azure ML}، یک پلتفرم بومی \lr{MLOps} طراحی و پیاده‌سازی شده است. این پلتفرم با استفاده از ابزارهای متن‌باز توسعه داده شده و قابلیت‌های مختلفی از جمله مدیریت چرخه عمر مدل، نظارت بر عملکرد، به‌روزرسانی خودکار و تعامل با ابزارهای موجود را فراهم می‌کند. این تحقیق با هدف پاسخ به سوالات تحقیقاتی مرتبط با نیازهای یک پلتفرم \lr{MLOps} مدرن و امکان‌پذیری پیاده‌سازی آن با استفاده از ابزارهای متن‌باز انجام شده است. نتایج حاصل از پیاده‌سازی و آزمایش‌ها نشان می‌دهند که پلتفرم پیشنهادی نه تنها به نیازهای مدرن پاسخ می‌دهد بلکه استقلال از ارائه‌دهندگان ابر را نیز فراهم می‌آورد. نهایتاً، یافته‌های تحقیق و پیشنهادات برای تحقیقات آینده ارائه شده‌اند.

\vspace{15mm}
\noindent \textbf{کلمات کلیدی:}
\vskip 2mm
\begin{tabular}{rr}

1-پیاده‌سازی یادگیری ماشین & \lr{Machine Learning Operation (MLOps).} \\
2-ادغام مداوم &‌ \lr{Continuous Integration (CI).} \\
3-استقرار مداوم &‌ \lr{Continuous Deployment (CD).} \\
4- مدیریت چرخه‌عمر یادگیری ماشین &‌ \lr{Machine Learning Lifecycle Management.} \\
5-پلتفرم ابری متن‌باز & \lr{Open-Source Cloud platform.} \\

\end{tabular}

%\noindent \rule{\textwidth}{2pt}


\end{minipage}
\end{center}
\vfil



% شروع فهرست مطالب و اشکال و غیره %%%%%%%%%%%%%%%%%%%%%%%%%%%%%%
\setcounter{page}{1}
\let\oldthepage=\thepage %Backup old definition
\renewcommand{\thepage}{\beginR\adadi{page}\endR}

\tableofcontents
\listoftables
\listoffigures
\chapter*{فهرست کلمات اختصاری}  
%\chapter*{\vspace*{-2cm} فهرست کلمات اختصاری}  
% دستور \vspace*{-4cm} برای این گذاشته شده است که  در صورت طولانی‌تر شدن این جدول از آن (با عدد مناسب) استفاده
% شود تا شروع صفحه از بالاتر باشد و در نهایت از شکسته‌شدن صفحه  جلوگیری شود (چون tabular قابلیت شکستن 
% صفحه ندارد). 
% توجه: یک روش دستی برای داشتن یک جدول چندصفحه‌ای آن است که هر عنصر جدول زیر را جداگانه 
% در یک جدول نوشت و پهنای ستون اول را با p{?cm} تعیین کرد.



\begin{latin}
{
\setlength{\baselineskip}{0.5\baselineskip}
\hfil
\begin{tabular}{p{3cm}l}
DevOps & Development and Operation \\
MLOps  & Machine Learning Operation \\
CI     & Continuous Integration \\
CD     &‌ Continuous Deployment \\
CT     &‌ Continuous Training \\
VM     & Virtual Machine \\
IaC    & Infrastructure as Code \\
ETL‌    & Extract, Transform, Load \\
PVC    &‌ Persistent Volume Claim \\
RBAC   & Role-Based Access Control \\
CNN    & Convolutional  Neural Network \\
LSTM‌   & Long Short-Term Memory \\


\end{tabular}
}
\end{latin}


%%%%%%%%%%%%%%%%%%%%%%%
\setcounter{page}{1}
\let\thepage=\oldthepage % Restore the definition 

% شروع فصلهای تز (متن اصلی تز) %%%%%%%%%%%%%%%%%%%%%%%%%%%%%%%%%%

\chapter{مقدمه} \label{ch:Introduction}

گسترش روز افزون شبکه جهانی اینترنت و توسعه فناوری اطلاعات، نیاز فزاینده‌ای را به استفاده از
سرویسهای چندرسانه‌ای دیجیتال،  در پی داشته به طوریکه کاربردهای دیجیتال شاهد رشد
شگرفی در طول دهه گذشته بوده است که نتیجه  آن ایجاد سیستمهای  کارآمد در ذخیره، انتقال 
و بازیابی اطلاعات است. مزایای فراوان فناوری دیجیتال، باعث محبوبیت و کاربرد هر~
چه بیشتر آن توسط اشخاص شده تا جاییکه حتی وسایل ضبط و پخش صدا و تصویر آنالوگ خانگی
هم به سرعت با نمونه‌های دیجیتال جایگزین شده‌اند. اما این موضوع مسائل حاشیه‌ای دیگری
برای بشر ایجاد نموده است. به طوریکه امکان تهیه کپی‌های متعدد از روی نسخه اصلی 
بدون کاهش کیفیت آن و یا سادگی جعل و تغییر محتوای اطلاعاتی نسخه اصلی، باعث شده که 
مالکیت  معنوی\footnote{\lr{Intellectual Property}} صاحبان اثر به خطر افتاده و در 
نتیجه بسیاری از ارائه دهندگان سرویسهای چندرسانه‌ای (از جمله شرکتهای فیلم‌سازی)
از ارائه نمونه دیجیتال محصولاتشان خودداری نمایند. لذا برطرف نمودن این مشکلات،
یکی از زمینه‌های پژوهشی مهم در عرصه مخابرات و بخصوص پردازش سیگنال است. 

              
    

\chapter{مروری بر جداسازی کور منابع (\lr{BSS}) و تجزیه به مؤلفه‌های مستقل (\lr{ICA})}

%%%%%%%%%%%%%%%%%%%%%%%%%%%%%%%%%%%%%%%%%%%%%%%%%%%%%%%%%%%%%%%%%%%%
%%%%%%%%%%%%%%%%%%%%%%%%%%%%%%%%%%%%%%%%%%%%%%%%%%%%%%%%%%%%%%%%%%%%
\section[جداسازی کور منابع (\lr{BSS}) و تجزیه به مؤلفه‌های مستقل (\lr{ICA})]{جداسازی 
کور منابع (\lr{BSS}\protect\footnote{\lr{Blind Source Separation}}) و تجزیه به مؤلفه‌های مستقل 
(\lr{ICA}\protect\footnote{\lr{Independent Component Analysis}})}

%%%%%%%%%%%%%%%%%%%%%%%%%%%%%%%%%%%%%%%%%%%%%%%%%%%%%%%%%%%%%%%%%%%%
\subsection{معرفی مسأله \lr{BSS}}
فرض کنید $N$  منبع صوتی  مختلف در اتاقی قرار گرفته‌اند و سیگنالهای
صوتی این منابع را  $s_1(t)$ تا  $s_N(t)$ بنامید. همچنین فرض کنید که در این 
اتاق  $N$ میکروفون قرار دارد و سیگنال دریافتی این میکروفونها را با 
$x_1(t)$ تا  $x_N(t)$ نشان دهید. در اینصورت هر میکروفون ترکیبی از کلیه 
سیگنالهای  $s_1(t)$ تا  $s_N(t)$ را دریافت می‌کند، یعنی:
\begin{equation}
\label{eq: First Lin Model}
x_i(t) = a_{i1} s_1(t) + a_{i2} s_2(t) + \cdots + a_{iN} s_N(t)
\end{equation}
که در آن ضرایب $a_{ij}$ نامعلوم هستند. همچنین توجه کنید که در اینجا 
ساده‌ترین 
حالت مسأله در نظر گرفته شده است، یعنی تعداد میکروفونها با تعداد منابع صوتی 
مساوی فرض شده و از تأخیرِ متفاوت منابع در رسیدن به میکروفونهای مختلف صرفنظر 
شده است. اکنون سؤال اینجاست که آیا می‌توان تنها با در دست داشتن خروجی 
میکروفونها ($x_1(t)$ تا $x_N(t)$)، سیگنالهای منابع ($s_1(t)$ تا $s_N(t)$) 
را بدست آورد؟ توجه شود که نه هیچگونه اطلاعاتی در مورد منابعِ $s_1(t)$ تا 
$s_N(t)$ در دست داریم و نه هیچ اطلاعی در مورد نحوه مخلوط‌شدن آنها 
(ضرایب  $a_{ij}$).

ممکن است در وهله اول، این مساله  \lr{ill posed} به نظر برسد و اطلاعات موجود
برای حل مساله کافی به نظر نیاید. اما توجه می‌کنیم که مغز انسان این
جداسازی را به راحتی انجام می‌دهد: وقتی در یک مهمانی شلوغ با یک نفر
صحبت می‌کنید، به راحتی صدای وی را از سایر صداهای موجود در محیط جدا
می‌کنید، حتی اگر سروصدای بقیه از صدای مخاطب شما بلندتر باشد (به همین
دلیل، مساله جداسازی چند سیگنالِ صحبتِ مخلوط، در متون علمی
\lr{Cocktail party problem} نامیده می‌شود).

اگرچه در اینجا ما مسأله \lr{BSS}  را برای سیگنالهای صوتی بیان کردیم، ولی مدل
معادلهٔ (\ref{eq: First Lin Model}) می‌تواند در کاربردهای دیگری نیز ظاهر شود، 
مثلاً در جداسازی سیگنال  \lr{ECG} جنین از مادر، جداسازی \lr{artifact}ها (مثل اثر پلک‌زدن) 
از سیگنال مغزی \lr{EEG}، کاربرد در مخابرات سیار، جداسازی تصاویر، استخراج ویژگی‌ها
(\lr{Feature extraction}) در پردازش تصاویر و غیره. در حالت کلی این مساله با
عنوان \lr{Blind Source Separation}  (\lr{BSS}) شناخته می‌شود (کلمه  \lr{Blind} با اشاره
به این موضوع است که نه هیچ اطلاعی در مورد سیگنالهای مخلوط~شده داریم و
نه در مورد نحوه مخلوط~شدن آنها).

%%%%%%%%%%%%%%%%%%%%%%%%%%%%%%%%%%%%%%%%%%%%%%%%%%%%%%%%%%%%%%%%%%%%
\subsection{تاریخه \lr{BSS}}
با توجه به ظاهر  \lr{ill posed} مسألهٔ \lr{BSS}، مطرح شدن آن در پردازش سیگنال نسبتاً
دیرهنگام بوده است. امکان حل مسألهٔ \lr{BSS} نخستین بار در سال  ۱۹۸۵ و توسط
هِرو (\lr{H\'erault})، ژوتن (\lr{Jutten}) و اَنس (\lr{Ans}) در یک کاربرد بیولوژیکی ارائه شد 
\cite{HeraJA85} (اولین مقاله بین‌المللی روی این موضوع به زبان انگلیسی
\cite{HeraJ86} است). ایده اولیه حدود سال ۱۹۸۲ و هنگامی بوجود آمد
که این تیم روی یک مساله بیولوژیکی کار می‌کرد و مشاهده کردند که هنگام
انتقال اطلاعاتِ مربوط به سرعت و مکان زاویه‌ایِ یک مفصل به مغز، اطلاعات
مربوط به سرعت و مکان توسط دو سیگنال عصبی مختلف به مغز ارسال
می‌شوند، ولی در هر دوی این سیگنالها اطلاعات مکان و سرعت با هم مخلوط
هستند و در نتیجه مغز حتماً باید روشی برای جداکردن این اطلاعات داشته باشد.
به عبارتی مغز باید به طریقی بتواند سیستم معکوس را تنها از روی مشاهدات
تخمین بزند. آنها در نهایت روشی برای تخمین پارامترهای سیستم معکوس یافتند
که بعدها به آلگوریتم \lr{HJ} (\lr{H\'erault-Jutten}) معروف شد.

    \begin{figure}[tb]
     \centering
     \unitlength 0.5cm \linethickness{0.4pt}
     \begin{LTR}
     \begin{picture}(8.25,4)(-1,0)
       \put(0,0){\vector(1,0){2}}
       \put(0,3){\vector(1,0){2}}

       \put(-1,-0.5){\makebox(1,1)[cc]{$x_2$}}
       \put(-1,2.5){\makebox(1,1)[cc]{$x_1$}}

       \put(2.25,0){\circle{0.5}}
       \put(2.25,3){\circle{0.5}}
       \put(2.1,0){\line(1,0){0.3}}
       \put(2.1,3){\line(1,0){0.3}}
       \put(2.25,-0.15){\line(0,1){0.3}}
       \put(2.25,2.85){\line(0,1){0.3}}

       \put(3.2,2.3){\makebox(1,1)[cb]{{\small $-m_{12}$}}}
       \put(3.2,0.3){\makebox(1,1)[cb]{{\small $-m_{21}$}}}


       \put(2.5,0){\vector(1,0){4}}
       \put(2.5,3){\vector(1,0){4}}
       \put(6.5,-0.5){\makebox(1,1)[cc]{$y_2$}}
       \put(6.5,2.5){\makebox(1,1)[cc]{$y_1$}}

       \put(5.25,0){\vector(-1,1){2}}
       \put(5.25,3){\vector(-1,-1){2}}

       \put(5.25,0){\line(-1,1){2.82}}
       \put(5.25,3){\line(-1,-1){2.82}}

       \put(5.25,0){\circle*{0.3}}
       \put(5.25,3){\circle*{0.3}}
     \end{picture}
     \end{LTR}     

\label{fig: HJ Sep Sys}
%\farsi 
\caption{سیستم جداساز در روش \lr{HJ}}
    \end{figure}
هِرو و  ژوتن برای جداسازی دو سیگنال مخلوط‍، ابتدا مدل فیدبک‌دار شکل~\ref{fig: HJ Sep Sys}
را برای سیستم جداساز در نظر گرفتند.
روشن است که این مدل با انتخاب مناسب ضرایب  $m_{12}$ و $m_{21}$ قادر است 
سیگنالها را از هم جدا کند. در واقع داریم:
\begin{equation}
    \left\{\begin{array}{rcc}
        y_1 &\eqnarrayeq& x_1 - m_{12} y_2 \\
        y_2 &\eqnarrayeq& x_2 - m_{21} y_1
      \end{array}
    \right.
\end{equation}
و یا به فرم ماتریسی  $\yb=\xb-\Mb\yb$، که در آن:
\begin{equation}
\Mb \triangleq \left[\begin{array}{cc} 1 & m_{12} \\
m_{21} & 2
\end{array}
\right]
\end{equation}
در نتیجه خواهیم داشت:
\begin{equation}
\yb = (\Ib + \Mb)^{-1} \xb
\end{equation}
بنابراین اگر ضرایب $m_{12}$ و $m_{21}$ بگونه‌ای انتخاب شوند که $\Ib+\Mb=\Ab$ 
شود (که در آن  $A\triangleq[a_{ij}]$ ماتریس ضرایب مخلوط‌کننده است)، 
منابع از هم جدا خواهند شد. البته این معادلات تنها نشان می‌دهند که سیستم 
شکل~\ref{fig: HJ Sep Sys} قادر به جداسازی منابع است و هیچ راه عملیی برای تعیین
پارامترهای سیستم جداساز بدست نمی‌دهد، چراکه ماتریس مخلوط  $\Ab$
مجهول است. سپس آلگوریتم \lr{HJ} (به عنوان اولین آلگوریتم جداسازی منابع و اولین
آلگوریتمی که قابل حل بودن مساله \lr{BSS} را نشان داد) بر مبنای این ایده استوار 
است که با توجه به اینکه منابع از هم مستقل آماری هستند، ضرایب  
$m_{12}$ و $m_{21}$ باید بگونه‌ای تعیین شوند که خروجی‌های $y_1$ و $y_2$ 
هم مستقل آماری باشند. از طرف دیگر دو تابع غیرخطی و فرد $f$ و $g$ را در 
نظر بگیرید .اگر $E\{f(y_1)g(y_2)\}$ و $E\{f(y_2)g(y_1)\}$  برابر صفر شوند
(میانگین سیگنالها صفر فرض شده است)، از روی بسط تیلور این توابع می‌توان انتظار
داشت که کلیه ممانهای متقابل (از مرتبه فرد $y_1$ و $y_2$) صفر باشند که به 
معنی استقلال آماری آنهاست. در نهایت  هِرو و ژوتن آلگوریتم زیر را برای تعیین ضرایب 
$m_{12}$ و  $m_{21}$ پیشنهاد کردند:
\begin{equation}
    \left\{
      \begin{array}{rcc}
        m_{12} \leftarrow m_{12} - \mu f(y_1) g(y_2) \\
        m_{21} \leftarrow m_{21} - \mu f(y_2) g(y_1)
      \end{array}
    \right.
\end{equation}
که در آن  عدد مثبت کوچکی است که  \lr{learning rate} را تعیین می‌کند. سپس  هِرو و 
ژوتن توابع فرد ساده‌ای مثل $f(y)=y^3$ و $g(y)=\arctan(y)$ را در نظر گرفتند 
و نشان دادند که این آلگوریتم واقعاً قادر است سیگنالهای مخلوط را جدا کند، به عبارتی حل 
مسألهٔ \lr{BSS} ممکن است. 

زیبایی و سادگی آلگوریتم \lr{HJ} در حل یک مسألهٔ ظاهراً پیچیده به تدریج توجه
محققان دیگر را نیز به این مسأله جلب کرد. برخی از آنها (مثل  \lr{Lacoume}) در
ابتدا تنها برای آنکه اشتباه بودن ادعای  \lr{HJ} را نشان دهند به کار روی این مسأله
پرداختند، که این تلاشها در نهایت منجر به اثبات ریاضی قابل حل بودن \lr{BSS}، و
ارائه آلگوریتم‌های بهتری برای آن شد \cite{Card89b,LacoR88,Como89c}.
با این وجود تا میانه دهه ۱۹۹۰، این مسأله بیشتر در بین محققان فرانسوی باقی
ماند و معدود مقالات بین‌المللی ارائه شده در کنفرانسهای شبکه‌های عصبی، در
بین انبوه مقالات روی موضوعاتی چون \lr{back-propagation}، شبکه‌های هاپفیلد و
\lr{Kohonen's self organizing maps} که در آن موقع مورد توجه زیاد محققان
شبکه‌های عصبی بودند، مدفون شد.

زمینه متفاوت دیگری که در آن موقع در حال گسترش بود،  
\lr{Higher Order Spectral Analysis} (\lr{HOS}) بود که نخستین کنفرانس آن در سال  
۱۹۸۹ برگزار شد و کارهای \lr{Cardoso} و  \lr{Comon} در این کنفرانس ارائه شدند
\cite{Card89b,Como89c}.

از حدود میانه دهه ۱۹۹۰ و بعد از آنکه \lr{Bell} و \lr{Sejnowski} روش خود را بر مبنای
\lr{Infomax} ارائه کردند \cite{BellS95}، این مسأله مورد توجه بیشتر
محققان قرار گرفت و آلگوریتم‌های مختلفی برای حل مسأله در حالتهای مختلف
ارائه گردید. بخصوص باید به روش \cite{CardL96} اشاره کرد که خانواده‌ای از 
روشهای \lr{equivariant} برای جداسازی منابع معرفی می‌کند. در این روشها، 
کیفیت جداسازی به  \lr{condition number} ماتریس مخلوط‌کننده بستگی
ندارد، یعنی هرچه هم که سیگنالها به شدت مخلوط شده باشند (خروجی
سنسورها به هم شبیه باشند، یا دترمینان ماتریس مخلوط‌کننده به صفر نزدیک باشد)، 
تأثیری در کیفیت نهایی جواب ندارند (در عدم حضور نویز).

از حدود میانه دهه ۹۰ مقالات در این زمینه به سرعت افزایش یافت. اولین
کنفرانس بین‌المللی روی این موضوع (\lr{ICA99}) در سال ۹۹  در فرانسه برگزار شد 
که در آن بیش از یکصد محقق که روی این مساله کار می‌کردند، 
گرد هم آمدند. این کنفرانس از آن سال تا کنون هر $1.5$ سال یکبار و
در کشورهای مختلف برگزار می‌شود (\lr{ICA2009} قرار است در برزیل برگزار شود).

%%%%%%%%%%%%%%%%%%%%%%%%%%%%%%%%%%%%%%%%%%
\subsection{کاربردهای \lr{BSS}}
با آنکه مساله \lr{ICA} یا \lr{BSS} نسبتاً مبحث جدیدی در پردازش سیگنال است، تا
کنون استفاده از آن در کاربردهای مختلفی گزارش شده است. از جمله این
کاربردها می‌توان موارد زیر را برشمرد:
\begin{itemize}
\item کاربرد در  \lr{Feature extraction} در پردازش تصویر.

\item جداسازی تصاویر مخلوط.

\item کاربرد در \lr{Denosing} سیگنال تصویر.

\item استفاده در جداسازی سیگنالهای پزشکی \lr{EEG}،  \lr{ECG} و \lr{MEG}.

\item استفاده در حذف سروصدای ماشین از صدای ضبط شده در گوشی تلفن همراه.

\item استفاده در حذف نویز و سروصدا به عنوان یک بلوک پیش~پردازش در
سیستمهای تشخیص صحبت (\lr{Speech recognition}). با استفاده از چنین بلوکی می‌توان 
نرخ تشخیص صحبت را افزایش داد.

\item حل  \lr{cocktail party problem} (جداسازی چند سیگنال صحبت مخلوط). این
کاربرد یکی از مشکلترین کاربردهای  \lr{BSS} است که با آنکه برای مخلوطهای
مصنوعی و سیگنالهایی که در یک اتاق آکوستیک ضبط شده‌اند، بخوبی کار
می‌کند، برای جداسازی مخلوط ضبط~شده در یک اتاق معمولی (وقتی از
روی دیوارها اکو وجود دارد) هنوز کیفیت روشهای ارائه‌شده چندان بالا
نیست.

\item استفاده در تشخیص همزمان چند جسم در سیستم‌های تشخیص شیء 
با استفاده از القای الکترومغناطیسی. در این سیستم‌ها هر شخص یا
شیء یک  \lr{Label} دارد که حضور آن در نزدیکی سنسور مربوطه باعث
شناسایی شخص یا شیء می‌شود. با استفاده از  \lr{ICA} امکان
\lr{identification} همزمان چند  \lr{label} بوجود می‌آید (وقتی چند  \lr{label} در
نزدیکی دستگاه تشخیص ظاهر می‌شوند). این کاربرد  \lr{patent} شده است.

\item استفاده در مخابرات (\lr{CDMA}).

\item کاربرد در اقتصاد (استخراج عوامل پنهان در اطلاعات تجاری).

\item کاربرد در شناسایی جو: یک سیگنال به طرف جو فرستاده می‌شود،
انعکاسهای آن از روی لایه‌های مختلف جو با استفاده از  \lr{ICA} از هم جدا
می‌شوند که هرکدام در برگیرنده اطلاعاتی در مورد لایه‌های مختلف جو
است.

\item کاربرد در زمین‌شناسی: مشابه روش قبل برای شناسایی لایه‌های مختلف زمین.

\item کاربرد در ستاره‌شناسی: مؤلفه‌های مستقل تصاویر دریافتی از فضا با استفاده 
از  \lr{ICA} استخراج می‌شوند و از روی آنها اطلاعاتی در مورد منابع
مستقلی که آنها را ارسال کرده‌اند بدست می‌آید.
\end{itemize}

%%%%%%%%%%%%%%%%%%%%%%%%%%%%%%%%%%%%%
\subsection{\lr{ICA} به عنوان روشی برای حل مسأله \lr{BSS}}
دیدیم که ایده اصلی حل مسأله \lr{BSS}
استفاده از  استقلال منابع و بیرون کشیدن مؤلفه‌های مستقلِ
مشاهدات است. به همین دلیل، گرفتن یک بردار مشاهده و تجزیه آن به مؤلفه‌های مستقل را
«\underline{تجزیه به مولفه‌های مستقل}» 
یا\lr{ICA} (\lr{Independent Component Analysis}) نام‌گذاری کردند. به عبارت دیگر،
\lr{ICA} به هر آلگوریتم یا روشی اطلاق می‌شود که نمونه‌هایی از یک بردار تصادفی را به عنوان ورودی  
دریافت کرده، و در خروجی مؤلفه‌های مستقل آن را بدست می‌دهد، یعنی مؤلفه‌هایی می‌دهد که از یکدیگر
بطور آماری  مستقل هستند و بردار ورودی را می‌توان به صورت ترکیب این مؤلفه‌های مستقل بیان کرد.

این نام‌گذاری در مقایسه با تبدیل قدیمی‌تر «تجزیه به مؤلفه‌های اصلی 
(\lr{PCA}\footnote{\lr{Principal Component Analysis}})» انجام گرفته است: \lr{PCA} نمونه‌های یک بردار تصادفی را
دریافت کرده و تحت یک تبدیل ماتریسی دوران، خروجیهایی بدست می‌دهد که با یکدیگر ناهمبسته (\lr{decorrelated})
هستند و ورودیها را می‌توان به صورت ترکیب خطی این مؤلفه‌های ناهمبسته بیان کرد.

بنابراین تفاوت \lr{PCA} و \lr{ICA} آن است که \lr{PCA} به تبدیلی از نوع دوران محدود است ولی \lr{ICA} نیست. در عوض،
\lr{PCA} خروجیهایی می‌دهد که از یکدیگر ناهمبسته هستند (اسقلال درجه ۲)، ولی لزوماً مستقل آماری نیستند، اما
\lr{ICA} خروجیهایی می‌دهد که به مفهوم دقیق مستقل (و نه فقط ناهمبسته) هستند. 
به عنوان مثال به سادگی می‌توان نشان داد که 
در \lr{BSS} بدست آوردن خروجی‌هایی که از هم ناهمبسته باشند برای بازیابی منابع کافی نیست و ناهمبستگی (به عنوان
تقریبی از استقلال) نمی‌تواند مسأله \lr{BSS} را حل کند، و برای حل این مسأله استقلال کامل نیاز است. به عبارتی
نمی‌توان \lr{PCA} را برای حل مسأله \lr{BSS} به کار برد.

البته توجه داریم که \lr{ICA} دو ابهام ذاتی نیز دارد: تغییر ترتیب مؤلفهای مستقل و نیز تغییر انرژی آنها
تغییری در استقلال آنها نمی‌دهد. مثلاً در \lr{BSS} که تنها اطلاعات ما از منابع استقلال آماری آنها است، 
منابع را می‌توان از روی مخلوطها به دست آورد ولی با ابهام در انرژی و ترتیب آنها؛ به عبارت دیگر،
انرژی و ترتیب منابع را نمی‌توان تنها از روی فرض استقلال آماری آنها به دست آورد. البته 
بازیابی منابع با این دو ابهام
در بسیاری از کاربردها قابل قبول است  و اهمیت چندانی ندارد.

در حالتیکه تعداد سنسورها مساوی با تعداد منابع باشد، \lr{ICA} متداول‌ترین روش برای \lr{BSS} است. البته تعمیم 
\lr{ICA} به حالتی که تعداد سنسورها از منابع بیشتر باشد نیز بسیار ساده است.
برای مطالعه بیشتر در مورد \lr{ICA} خواننده را به 
مقالات مروری~\cite{MansBO00,LeeGBS99,HyvaO00,Hyva99Survey,Card98}، 
کتابهای~\cite{CichA02,HyvaKO01} و مراجع ذکر~شده در آنها ارجاع می‌دهیم.



\chapter{\lr{MLOps}}\label{ch:HVS}
 
\section{مقدمه}
 
در حالی که مدل‌های یادگیری ماشین به طور گسترده‌ توسعه یافته‌اند، انتقال آن‌ها از مفهوم آزمایشی به محیط تولید اغلب با شکست مواجه می شود. این فاصله بیشتر به خاطر این است که تاکنون توجه اصلی روی ساخت مدل‌ها بوده است، نه روی تولید محصولات یادگیری ماشین که قابلیت استفاده در محیط تولید را دارند. علاوه بر آن، مدیریت بخش‌ها و زیرساخت‌های پیچیده‌ای که برای یک استقرار موثر ضروری هستند نیز در این امر مغفول مانده اند. برای رفع این مسئله، مفهوم عملیات یادگیری ماشین یا \lr{MLOps} معرفی شده است. \lr{MLOps} بر روی خودکارسازی و عملیاتی کردن فرآیندهای یادگیری ماشین تمرکز دارد تا انتقال پروژه‌های یادگیری ماشین از مفهوم به تولید را تسهیل کند. این رویکرد شامل دیدگاه جامعی از طراحی سیستم، هماهنگی اجزا، تعریف نقش‌ها و مسئولیت‌ها می باشد. هدف کاهش خطا به منظور افزایش قابلیت اطمینان و کارایی سیستم‌های یادگیری ماشین در کاربردهای واقعی می باشد. این فصل به بررسی تعریف، اصول، ابزار و معماری جامعی از یک پلتفرم  \lr{MLOps} پرداخته و در نهایت، محصولات و رقبا این حوزه را بررسی می کنیم.

\section{تعریف مفاهیم اولیه}

\lr{MLOps}
 یا عملیات یادگیری ماشین به مجموعه‌ای از فرایندها، ابزارها و شیوه‌ها جهت مدیریت چرخه توسعه مدل‌های یادگیری ماشین در یک محیط عملیاتی اشاره دارد. همچنین این چرخه شامل همکاری بین دانشمندان داده و مهندسان \lr{DevOps} است به گونه ای که این اطمینان حاصل شود که مدل‌ها به طور مؤثر توسعه، استقرار، پایش و به‌روزرسانی می‌شوند. هدف \lr{MLOps} افزایش سرعت، قابلیت اطمینان و مقیاس پذیری مدل‌های یادگیری ماشین و فرایند توسعه این مدل ها در تولید است؛ درحالی‌که خطرات ناشی از ریسک عدم موفقیت را نیز کاهش می‌دهد. همچنین به‌کارگیری \lr{MLOps} فرایند مدیریت را ساده‌تر کرده، کیفیت را افزایش می‌دهد و استقرار مدل‌های یادگیری عمیق و یادگیری ماشین در محیط‌های تولید با مقیاس بزرگ را خودکار می‌کند. لذا می توان گفت یکی از اهداف \lr{MLOps}، بهبود خودکارسازی و ارتقای کیفیت مدل‌های تولید و درعین‌حال توجه به الزامات تجاری و نظارتی است. 
 
 
 استقرار مدل‌های یادگیری ماشین روی محیط عملیاتی در \lr{MLOps} اهمیت زیادی دارد، زیرا به سازمان‌ها کمک می‌کند تا مطمئن شوند که مدل‌هایشان در طول زمان دقیق، قابل‌اعتماد و کارآمد هستند. به‌طورکلی، \lr{MLOps} با خودکار کردن بسیاری از مراحل مربوط به استقرار و مدیریت مدل‌های یادگیری ماشین، به دانشمندان و مهندسان داده اجازه می‌دهد تا با همکاری یکدیگر به ارائه سریع‌تر و کارآمدتر مدل‌های یادگیری ماشین دست یابند. 
\subsection{اصول}
 برای تسهیل در رسیدن به اهداف فوق، تیم‌های \lr{MLOps} از اصول زیر استفاده می کنند:
\begin{itemize}
	\item 
	خط لوله خودکار \lr{CI/CD} و هماهنگ سازی جریان کاری
	\item 
	کنترل نسخه مدل‌های یادگیری ماشین، مجموعه‌داده‌ها و کد منبع
	\item 
	آموزش و ارزیابی مدوام مدل یادگیری ماشین
	\item 
	نظارت مستمر و ثبت گزارش برای ردیابی عملکرد مدل
	\item 
	حلقه های بازخورد\footnote{\lr{feedback loops}}
\end{itemize}




\chapter{آنتروپی و استفاده از آن در نهان~نگاری}\label{ch:Entropy}
 
\section{مقدمه}

در فصل گذشته سیستم بینایی انسان و ویژگی‌های آن را مورد بررسی قرار دادیم.
در حوزهٔ تبدیل \lr{DCT} و تبدیل موجک، دو مدل معروف و موجود برای سیستم بینایی
معرفی کردیم. همچنین  چارچوب کلی استفاده از مدلهای بینایی را برای  کاربرد
نهان نگاری، مطرح نمودیم و دو طرح مرجع \lr{P\&Z } و \lr{K\&R} را نیز مورد بررسی قرار دادیم. 
در این فصل به بیان هدف اصلی این پایان نامه که بررسی اثر آنتروپی در نهان~نگاری است، می‌پردازیم.
لذا ابتدا در بخش \ref{sec:entropy} به  بیان مفهوم آنتروپی پرداخته، سپس در 
بخش \dots 

\section{آنتروپی}\label{sec:entropy}

در این قسمت ابتدا توضیح مختصری در بارهٔ مفهوم آنتروپی داده خواهد شد. 


\chapter{پیاده سازی و نتایج}

\section{پیاده سازی پلتفرم}
\subsection{سیستم مدیریت}

در پلتفرم‌های بزرگ، نیاز به سیستمی برای مدیریت پلتفرم به وضوح احساس می‌شود. این سیستم باید قابلیت هماهنگی و یکپارچگی بین اجزای مختلف را داشته باشد تا اطمینان حاصل شود که همه بخش‌ها به درستی و بدون مشکل عمل می‌کنند. ویژگی‌های حیاتی این سیستم شامل استفاده از ابزارهای خودکارسازی و نظارت پیشرفته، مدیریت بهینه منابع سخت افزاری، پیاده‌سازی فرآیندهای مستمر بهبود و به‌روزرسانی، مدیریت دسترسی‌ها و امنیت و مستندسازی فرآیندها و تغییرات است. سیستم مدیریت شامل تمامی ابزارهای مدیریتی مانند مدیریت مخازن مولفه‌ها، کد و خط لوله \lr{CI/CD}، مانیتورینگ کل سیستم و جمع آوری لاگ است. علاوه بر این، استراتژی استقرار پروژه، اعمال مهاجرت‌ها\footnote{\lr{Migrations}}، پیکربندی پروژه‌ها، مدیریت سرورهای \lr{DNS} و \lr{NTP}\footnote{\lr{Network Time Protocol}} و استفاده از ابزارهایی مانند \lr{Foreman} برای نصب سیستم‌عامل به صورت \lr{PXE} نیز توسط همین سیستم مدیریت می‌شود.

به منظور پیاده سازی این سیستم، ما دو ماشین مجزا برای مدیریت پلتفرم به منظور ایجاد قابلیت تحمل خطا\footnote{\lr{Fault Tolerance}} و دسترسی پذیری بالا\footnote{\lr{High Availability}}  قرار می‌دهیم. از آنجایی که فرآیند و پروسه سنگینی روی این ماشین ها انجام نمی شود مشخصات کمتری می تواند به نسبت ماشین های پروژه داشته باشد. مشخصات سخت افزاری هرکدام از این ماشین ها در جدول
~\ref{tb: management conf}
قرار دارد.  این ماشین ها به صورت ماشین مجازی با استفاده از \lr{OpenStack} ساخته می شود. ماشین‌های مدیریت باید برای نصب و راه‌اندازی ابزارهای مدیریت پلتفرم پیکربندی شوند و این کار با استفاده از ابزار \lr{Ansible} انجام می گیرد. به این منظور \lr{Role} های مشخصی برای هر بخش نوشته شده است تا بتوان بدون هیچ کار دستی و به صورت کاملا خودکار سیستم ها را پیکربندی کرد. این \lr{Role} های انسیبلی با استفاده از قاعده نسخه گذاری \lr{Semantic} نسخه گذاری شده و در مخزن \lr{raw} موجود در \lr{Nexus} که یک ابزار مدیریت مخازن مولفه می باشد نگه داری خواهند شد. در نهایت برای پیکربندی سیستم، \lr{Role} موردنظر با نسخه مشخص از \lr{Nexus} گرفته شده و بااستفاده از \lr{Ansible} ماشین ها پیکربندی می شوند. از آنجایی که این پیکربندی در محیط های مختلف مانند توسعه و عملیات می تواند متفاوت باشد، ما با استفاده از قابلیت \lr{Overriding} در این ابزار مقادیر پیش فرض را برای هر محیط تغییر خواهیم داد. به همین منظور اسکریپتی طراحی شده که در لینک 
گیت هاب\footnote{\url{https://github.com/abolfazlyarian/mlops.git}} قابل مشاهده است.

\begin{table}
	\centering
	\caption{مشخصات سخت افزاری ماشین های مدیریت}
	\label{tb: management conf}
	\begin{tabular}{|c|c|c|c|}
		\hline
		\lr{OS} & \lr{Storage} &  \lr{RAM} & \lr{CPU} \\ \hline
		\lr{Ubuntu 18.04} & \lr{512 GB} & \lr{8 GB} & \lr{4 Core} \\ \hline
	\end{tabular}
\end{table}

پروسه پیکربندی و نصب ابزار در ماشین های مدیریت به صورت زیر انجام می گردد:
\begin{enumerate}
	\item 
	پیکربندی ماشین ها:
این قسمت شامل نصب و پیکربندی ابزارهایی نظیر
\lr{BIND}
برای سرور \lr{DNS}،
\lr{APT}
برای مدیریت ابزار در سیستم عامل \lr{Ubuntu}،
\lr{pip}
برای مدیریت کتابخانه ها پایتون،
\lr{chrony}
برای سرور \lr{NTP}
،\lr{LDAP} 
برای مدیریت کاربران و … می باشد.
	\item
نصب و پیکربندی \lr{Docker}: از آنجایی که مدیریت ابزارها به صورت کانتینر مناسب تر است، برای انجام مراحل بعدی نیاز به نصب \lr{Docker} می باشد. پس از نصب به منظور ذخیره سازی تمام مولفه ها مورد استفاده به مخزن ساخته شده در \lr{Nexus} مدیریت متصل خواهد شد.
	\item 
 \lr{Nexus}: از این ابزار به منظور مدیریت مخازن مولفه ها استفاده شده است. مخازن مورد استفاده ما \lr{APT}، 
	\lr{pip}،
	\lr{Docker}
	و \lr{raw} می باشد (شکل 
	~\ref{fig: nexus repo}).

	\item 
 \lr{Jenkins}: از این ابزار به منظور اجرا و مدیریت خط لوله های \lr{CI/CD} پروژها و هم چنین پیکربندی آن ها توسط مدیران سیستم استفاده می شود (شکل 
	~\ref{fig: jenkins}).
	
	\item 
 \lr{GitLab}: به منظور مدیریت کد در پروژه ها و هم چنین مدیریت \lr{Role}های انسیبلی برای پیکربندی پروژها استفاده می شود (شکل 
	~\ref{fig: gitlab}).
	
\begin{figure}[!t]
	\centering
	\includegraphics[scale=0.3]{nexus-repo.png}
	\caption{مخازن مولفه در \lr{Nexus}}
	\label{fig: nexus repo}
\end{figure}
\begin{figure}[!t]
	\centering
	\includegraphics[scale=0.3]{jenkins1.png}
	\caption{خط لوله ها \lr{CI/CD} در \lr{Jenkins}}
	\label{fig: jenkins}
\end{figure}
\begin{figure}[!t]
	\centering
	\includegraphics[scale=0.3]{gitlab.png}
	\caption{مخازن کد در \lr{Gitlab}}
	\label{fig: gitlab}
\end{figure}
\end{enumerate}
\clearpage
\subsection{خوشه کوبرنتیز}
در طراحی یک پلتفرم \lr{MLOps} جامع و کارآمد که تمامی ابزارهای مورد نیاز را در بر می‌گیرد، هدف اصلی ایجاد یک بستر یکپارچه، مقیاس‌پذیر و انعطاف‌پذیر برای مدیریت چرخه حیات مدل‌های یادگیری ماشین است. این پلتفرم شامل مجموعه‌ای از ابزارها و تکنولوژی‌های متن‌باز است که همگی روی خوشه کوبرنتیز مستقر می‌شوند. استفاده از کوبرنتیز در پلتفرم‌های \lr{MLOps} به دلیل قابلیت‌های منحصر به فرد آن در مدیریت خودکار، مقیاس‌پذیری و کارایی منابع است. 

برای نصب خوشه کوبرنتیز و انجام تست های اولیه پروژه \lr{MLOps}
، 4 ماشین مجازی با مشخصات ذکر شده در جدول 
~\ref{tb: mlops conf}
 استفاده شده است. در این پیاده سازی سه ماشین اول به عنوان گره اصلی و ماشین چهارم به عنوان گره کاری انتخاب شده است. همانند سیستم های مدیریت، ابتدا ماشین های مجازی که توسط \lr{OpenStack} ساخته خواهند شد پیکربندی می شوند. پس از آن با استفاده از ابزار \lr{Kubeadm} خوشه کوبرنتیز پیاده سازی خواهد شد.
\begin{table}
	\centering
	\caption{مشخصات سخت افزاری ماشین های خوشه کوبرنتیز}
	\label{tb: mlops conf}
	\begin{tabular}{|c|c|c|c|}
		\hline
		\lr{OS} & \lr{Storage} &  \lr{RAM} & \lr{CPU} \\ \hline
		\lr{Ubuntu 18.04} & \lr{2 TB} & \lr{128 GB} & \lr{40 Core} \\ \hline
	\end{tabular}
\end{table}

\subsubsection{پیاده سازی خوشه کوبرنتیز}
برای نصب و پیاده‌سازی کوبرنتیز با استفاده از \lr{kubeadm}، ابتدا پیش‌نیازهایی مانند \lr{Containerd} و بسته‌های کوبرنتیز شامل \lr{kubeadm}، \lr{kubelet} و \lr{kubectl} آماده می‌شوند. پس از پیکربندی سیستم و نصب \lr{Containerd}، بسته‌های مورد نیاز نصب شده و \lr{Swap} غیرفعال می‌گردد، چرا که کوبرنتیز برای عملکرد بهینه نیاز به غیرفعال بودن \lr{Swap} دارد. راه‌اندازی خوشه با اجرای دستور \lr{kubeadm init} بر روی گره اصلی آغاز می‌گردد. این دستور تنظیمات اولیه را انجام داده و فایل پیکربندی برای \lr{kubectl} ایجاد می‌کند که باید برای دسترسی به خوشه تنظیم گردد. پس از آن، نصب رابط شبکه کانتینر\footnote{\lr{Container Network Interface (CNI)}} انجام می‌شود و در این مورد، ابزار \lr{Calico} مورد استفاده قرار می‌گیرد. نصب \lr{Calico} امکان برقراری ارتباط بین پادها در داخل خوشه را فراهم می‌سازد.

در مرحله بعد، گره های کارگر به خوشه اضافه می‌شوند. دستورات لازم برای پیوستن نودهای کارگر به خوشه که از دستور قبل به‌دست آمده‌اند، بر روی هر گره کارگر اجرا می‌گردد تا این گره ها به گره اصلی متصل شوند. در نهایت، وضعیت گره ها و صحت پیوستن آن‌ها به خوشه با استفاده از \lr{kubectl get nodes} بررسی می‌گردد و اطمینان حاصل می‌شود که همه گره ها به درستی اضافه شده و آماده به کار هستند. این روش نصب، راهی سریع و مؤثر برای راه‌اندازی خوشه کوبرنتیز فراهم می‌آورد که امکان اضافه کردن گره  های جدید و مدیریت بهتر و مقیاس‌پذیری را به سهولت فراهم می‌سازد. تمامی این فرآیند با استفاده از \lr{Role} های انسیبلی روی ماشین های پروژه انچام می گردد.

\subsubsection{پیاده سازی مولفه های پروژه}
تمامی مولفه ها انتخاب شده در طراحی معماری فصل قبل با استفاده از \lr{Helm} که وظیفه مدیریت مولفه ها را دارد انجام می گیرد. \lr{Helm} یک ابزار قدرتمند برای مدیریت بسته‌های نرم‌افزاری کوبرنتیز است که فرآیند استقرار، بروزرسانی و مدیریت برنامه‌ها را ساده‌تر می‌کند. نصب یک چارت با استفاده از \lr{Helm} این امکان را می‌دهد که برنامه‌ها و سرویس‌ها با کمترین تلاش ممکن و به صورت اتوماتیک در خوشه کوبرنتیز مستقر شوند. چارت مورد نظر با استفاده از دستور \lr{helm install} به راحتی نصب می‌گردد. این دستور، چارت مورد نظر را از مخازن موجود در \lr{Nexus} دانلود کرده و آن را در خوشه کوبرنتیز مستقر می‌کند. این کار نیز با استفاده از یک \lr{Role} انسیبلی که نوشته شده است انجام می گیرد. 

یکی از مزایای استفاده از \lr{Helm}، قابلیت تنظیم پارامترهای چارت‌ها در زمان نصب است. این قابلیت امکان سفارشی‌سازی چارت‌ها را براساس نیازهای خاص پروژه فراهم می‌کند. با استفاده از این ویژگی و قابلیت مشابه در انسیبل، همانند قبل برای هر محیط مقادیر خاص آن محیط را پیکربندی می کنیم. علاوه بر این، \lr{Helm} با ارائه قابلیت بازگشت پذیری، امکان بازگشت به نسخه‌های قبلی چارت‌ها را نیز فراهم می‌سازد که این ویژگی برای مدیریت نسخه‌ها و رفع مشکلات بسیار مفید است. 

تمامی مولفه های اصلی نصب شده روی خوشه کوبرنتیز به منظور پیاده سازی پلتفرم \lr{MLOps} در شکل
??????????????
 نشان داده شده است. نکته ای که راجع به این مولفه ها لازم به ذکر است، به منظور افزایش قابلیت اطمینان و تحمل خطا، استفاده بیش از یک \lr{replica} برای هر مولفه ضروری است. این امر باعث می‌شود تا در صورت خرابی یکی از نمونه‌ها، مولفه مربوطه به کار خود ادامه دهد و سرویس‌دهی دچار اختلال نشود.
در شکل 
~\ref{fig: resource use no load}
نیز میزان مصرف منابع سخت افزاری سیستم زمانی که پلتفرم زیر بار نیست نشان داده شده است. همان طور که می بینید میزان مصرف منابع سخت افزاری به صورت یکنواخت به خوبی توزیع شده است که منجر به استفاده بهینه از منابع سخت افزاری می گردد.

\begin{figure}[!b]
	\centering
	\includegraphics[scale=0.7]{resource-use1.png}
	\caption{میزان مصرف منابع سخت افزاری پلتفرم بدون بار}
	\label{fig: resource use no load}
\end{figure}
\clearpage
\section{حل مسئله و نتایج}
 
پس از پیاده سازی پلتفرم \lr{MLOps} نیاز به ارزیابی آن داریم. هدف از این کار، بررسی قابلیت‌ها و عملکرد این پلتفرم در مدیریت و اجرای مدل‌های یادگیری ماشین به صورت خودکار و پیوسته است. به منظور ارزیابی و تست پلتفرم، سه مسئله مختلف مورد بررسی قرار می‌گیرند. در مسئله اول، از قابلیت \lr{Notebooks} در این پلتفرم استفاده شده است و دو مسئله دیگر به صورت کدهای پایتون نوشته شده‌اند که با بهره‌گیری از \lr{CI/CD}، روی پلتفرم پیاده‌سازی و اجرا می‌شوند. این رویکرد، به ما امکان می‌دهد تا عملکرد پلتفرم را در شرایط مختلف و با استفاده از ابزارهای متفاوت بررسی و ارزیابی کنیم.

\subsection{مسئله تشخیص ارقام}
\subsection{مسئله تحلیل احساسات در بازار سهام}

مسئله اصلی این پروژه تحلیل احساسات در بازار سهام است. هدف این است که با استفاده از نظرات و اخبار اقتصادی، احساسات مثبت و منفی موجود در بازار شناسایی شود و تأثیر آن بر رفتار بازار تحلیل گردد. داده‌های دادگان این مسئله به تجربیات احساس بازار سهام مربوط می‌شود  که از حساب‌های مختلف توئیتر، شامل 5791 توئیت مرتبط با اخبار اقتصادی پیرامون سهام، با دسته‌بندی احساس مثبت و منفی جمع‌آوری شده‌اند. این مجموعه‌داده به دو بخش مثبت و منفی تقسیم می شوند که تعداد موارد منفی 2106 و تعداد موارد مثبت 3685 است. در شکل 
~\ref{fig: stock data}
 ده سطر اول از این داده‌ها به عنوان نمونه نشان داده شده است. راه حل ارائه شده برای حل این مسئله استفاده از شبکه های \lr{LSTM} می باشد. به منظور کنترل و مدیریت این مسلئه از دو خط لوله مجزا برای آموزش و استقرار استفاده شده است. 
\begin{figure}[!b]
	\centering
	\includegraphics[scale=0.6]{stock-table.png}
	\caption{مجموعه داده احساسات در بازار سهام}
	\label{fig: stock data}
\end{figure}

\subsubsection{خط لوله آموزش}
در این خط لوله از قابلیت کانتینر برای استقرار آسان استفاده شده است و شامل آزمون‌هایی برای اطمینان از عملکرد و دقت راه‌حل‌های پیاده‌سازی شده است که هر کانتینر با ساخت و اجرای فایل و تصویر داکر مربوطه اجرا خواهد شد. هم چنین،‌ برای ساخت خط لوله در بستر پلتفرم لازم است که در ابتدا بخش‌های مختلف پروژه را به مولفه‌های مجزا تقسیم کرد، به طوری که هر مولفه دارای یک فایل داکر بوده و به صورت مجزا به آن نگاه خواهد شد. در این جا ما با سه مولفه سر و کار خواهیم داشت که مراحل پیش‌پردازش، آموزش و آزمایش چرخه‌ی یادگیری ماشین ما خواهند بود. همچنین برای ساخت خط لوله نیاز است تا اسکریپت جداگانه‌ای مربوط با کدهای این بخش که با کتابخانه \lr{kfp} است نوشته شود، تا بتوان با استفاده از تصاویر داکری ساخته شده در هر مولفه، پایپلاین مورد نظر را روی بستر پلتفرم \lr{MLOps} ساخت. در شکل 
~\ref{fig: stock train}
یکی از اجراهای خط لوله \lr{CI/CD} برای مرحله آموزش را مشاهده می کنید.
\begin{figure}[!t]
	\centering
	\includegraphics[scale=0.4]{stock-train.png}
	\caption{خط لوله آموزش مسئله تحلیل احساسات در بازار سهام}
	\label{fig: stock train}
\end{figure}

همان طور که می بینید این خط لوله شامل چندین بخش از جمله پیش پردازش،‌ آموزش،‌ ارزیابی و پیاده سازی روی بستر پلتفرم \lr{MLOps} با استفاده از \lr{Kubeflow Pipeline} می باشد.


پیش پردازش شامل چندین مرحله کلیدی است که داده‌ها را برای مدل‌سازی آماده می‌کند. این مراحل شامل تغییر برچسب‌ها به مقادیر دسته‌ای، تمیز کردن متن از کاراکترها و کلمات بی‌معنی، تقسیم داده‌ها به دو بخش آموزش و ارزیابی، و تبدیل متن به دنباله‌هایی از کلمات است.  همچنین، واژگان با استفاده از تعبیه‌گرهای پیش‌آموزش داده شده \lr{GloVe} به بردارهایی تبدیل می‌شوند که حاوی اطلاعات معنایی و ارتباطی کلمات هستند. فایل‌های تعبیه‌گر \lr{GloVe} از \lr{MinIO} دانلود شده و بردارهای تعبیه برای واژگان موجود در داده‌های آموزش ساخته می‌شود. این بردارها سپس برای ایجاد یک ماتریس تعبیه جهت استفاده در شبکه عصبی به کار می‌روند. داده‌های تعبیه‌شده و سایر داده‌های پردازش‌شده نیز در \lr{MinIO} ذخیره می‌شوند تا در مراحل بعدی مورد استفاده قرار گیرند. برای ذخیره و بازیابی داده‌ها، از دو ابزار اصلی \lr{MinIO} و \lr{PostgreSQL} استفاده می‌شود. در کلاس پیش‌پردازش، فایل‌های تعبیه‌گر \lr{GloVe} را از \lr{MinIO} دانلود و بردارهای تعبیه برای واژگان موجود در داده‌های آموزش ساخته می‌شود. همچنین، داده‌های پردازش‌شده نیز در \lr{MinIO} ذخیره می‌شوند. \lr{PostgreSQL} نیز برای بازیابی مجموعه‌داده اولیه به کار می‌رود، به طوری که داده‌ها از پایگاه‌داده خوانده شده و برای پردازش و آماده‌سازی بیشتر به کلاس پیش‌پردازش منتقل می‌شوند. 


در بخش آموزش از شبکه‌های حافظه کوتاه-مدت بلند\footnote{\lr{Long short-term memory (LSTM)}} به دلیل توانایی‌شان در درک و پردازش توالی‌های زمانی طولانی‌مدت، برای تحلیل احساسات متنی استفاده شده است. برخلاف مدل‌های سنتی که ممکن است نتوانند ارتباطات طولانی‌مدت بین کلمات را به خوبی درک کنند، \lr{LSTM} با استفاده از سلول‌های حافظه و مکانیزم‌های دروازه‌ای خود، قادر است اطلاعات مهم را حفظ و اطلاعات غیرضروری را فراموش کند. برای استفاده از \lr{LSTM} در تحلیل احساسات، از داده های پیش پردازش شده که در \lr{MinIO} ذخیره‌سازی شده بود استفاده می شود. پس از آماده‌سازی داده‌ها، مدل \lr{LSTM} با تعریف لایه‌های مختلف از جمله لایه‌های تعبیه‌گر، لایه‌های \lr{LSTM} دوطرفه و لایه‌های چگال ساخته می‌شود. این لایه‌ها به مدل امکان می‌دهند تا وابستگی‌های معنایی و زمانی بین کلمات را بهتر درک کند. علاوه بر این،‌ به منظور جلوگیری از بیش برازش از لایه های \lr{Dropout} با نرخ \lr{0.2} استفاده شده است.

در انتها پس از آموزش، مدل با استفاده از مجموعه دادگان ارزیابی، بررسی می شود. معیارهای ارزیابی مدل نیز دقت آن می باشد. در صورتی که بررسی های انجام شده از دقت مدل های گذشته بیشتر بود، مدل در \lr{MinIO} ذخیره‌سازی می شود.

به منظور اجرا چرخه یادگیری ماشین گفته شده این چرخه با استفاده از کتابخانه \lr{kfp} روی بستر پلتفرم اعمال شده تا فرآیند پیش پردازش،‌ آموزش و ارزیابی مدل را روی پلتفرم انجام دهد (شکل ~\ref{fig: stock workflow}).

\begin{figure}[t]
	\centering
	\includegraphics[scale=0.8]{stock-workflow.png}
	\caption{چرخه کاری آموزش مسئله تحلیل احساسات در بازار سهام}
	\label{fig: stock workflow}
\end{figure}

\subsubsection{خط لوله استقرار}
همان طور که در شکل 
~\ref{fig: stock serve}
مشاهده می کنید، اجرای خط لوله استنتاج مدل در پلتفرم با استفاده از \lr{Kserve} شامل چند مرحله کلیدی است. ابتدا، یک فایل پیکربندی برای \lr{Kserve} ایجاد می‌شود که جزئیات مدل، مسیر مدل ذخیره شده و هرگونه مراحل پیش‌پردازش یا پس‌پردازش لازم را مشخص می‌کند. سپس، یک جزء\footnote{\lr{Component}}  \lr{Kserve} با استفاده از کتابخانه \lr{kfp} ایجاد می‌شود که از فایل پیکربندی \lr{YAML} استفاده می‌کند. این جزء باید شامل منطق لازم برای استقرار مدل، مدیریت درخواست‌های ورودی و ارائه پیش‌بینی‌ها باشد. پس از ایجاد جزء، با استفاده از \lr{Kubeflow Piepline} روی بستر پلتفرم پیاده‌سازی می‌شود. این شامل تعریف وابستگی‌ها بین اجزا است تا اطمینان حاصل شود که جزء استقرار مدل پس از اجزا آموزش مدل و ارزیابی اجرا می‌شود. در نهایت، ارتباط با \lr{Endpoint} از طریق \lr{REST API} انجام می‌شود. این فرآیند به تحلیلگران اجازه می‌دهد تا مدل‌های یادگیری ماشین را به طور مؤثر استقرار و مدیریت کنند و از ویژگی‌های پیشرفته مانند مقیاس‌پذیری خودکار، تعادل بار و پایش عملکرد بهره‌مند شوند. 
\begin{figure}[!t]
	\centering
	\includegraphics[scale=0.4]{stock-serve.png}
	\caption{خط لوله استقرار مسئله تحلیل احساسات در بازار سهام}
	\label{fig: stock serve}
\end{figure}


\subsubsection{نتایج}
از آنجا که دقت و کارایی مدل هدف پیاده‌سازی مسئله نبوده و صرفا هدف ما تست عملکردی پلتفرم می باشد، ملاک ارزیابی ما زمان اجرا خط لوله ها و چرخه یادگیری ماشین،‌ مدت زمان پاسخ و میزان مصرف منابع می باشد. 
در جدول 
~\ref{tb: stock pipeline time}
زمان اجرای هر بخش را برای اجرای یک تغییر نشان داده شده است. کل مدت زمان آموزش تا استقرار مدل مجموعا 392 ثانیه می باشد.
توجه داشته باشید که این زمان بدون لحاط زمان کش می باشد. در صورتی که تغییر در یک جزء صورت می گیرد تنها آن جزء اجرا شده و بقیه اجزا از حافظه نهان برای اجرا آن استفاده می کنند. 
\begin{table}
	\centering
	\caption{زمان اجرا خط لوله ها \lr{CI/CD} و چرخه یادگیری ماشین در مسئله تحلیل احساسات در بازار سهام }
	\label{tb: stock pipeline time}
	\begin{tabular}{|c|c|}
		\hline
		زمان (ثانیه) & خط‌لوله  \\ \hline
		\lr{70} &  \lr{CI/CD} آموزش  \\ \hline
		\lr{187} &  چرخه کاری آموزش  \\ \hline
		\lr{37} &  \lr{CI/CD} استقرار  \\ \hline
		\lr{98} & چرخه کاری استقرار  \\ \hline
	\end{tabular}
\end{table}

به منظور تست مدل استقرار یافته،‌ به صورت همزمان به آن درخواست داده و زمان پاسخ و میزان مصرف \lr{CPU} را گزارش کرده ایم. نتایج این تست در جدول 
~\ref{tb: stock req}
گزارش شده است. 

\begin{table}
	\centering
	\caption{تست استنتاج مدل مسئله تحلیل احساسات در بازار سهام}
	\label{tb: stock req}
	\begin{tabular}{|c|c|c|}
		\hline
		\lr{CPU (mili cores)} &  زمان پاسخ (ثانیه) & تعداد  \\ \hline
		\lr{24}  & \lr{4.16} & \lr{1} \\ \hline
		\lr{110}  & \lr{4.59} & \lr{4} \\ \hline
		\lr{290}  & \lr{5.51} & \lr{10} \\ \hline
		\lr{387}  & \lr{6.48} & \lr{16} \\ \hline
		\lr{566}  & \lr{8.91} & \lr{32} \\ \hline
		\lr{845} & \lr{13.99} & \lr{64} \\ \hline
		\lr{1142}  & \lr{11.86} & \lr{100} \\ \hline
		\lr{1338}  & \lr{14.09} & \lr{128} \\ \hline
		\lr{1689}  & \lr{14.45} & \lr{256} \\ \hline
		\lr{2090} & \lr{15.8} & \lr{512} \\ \hline
	\end{tabular}
\end{table}

نتایج در تعداد درخواست ها بالا به دلیل استفاده از تکنیک مقیاس پذیری خودکار در پلتفرم با استفاده از ابزارهای \lr{Istio} و \lr{Knative} مناسب است ولی در تعداد درخواست های پایین مناسب نمی باشد. به منظور رفع این مشکل از سرویس \lr{LoadBalancer} در خوشه کوبرنتیز استفاده شده است که نتایج آن را در شکل 
~\ref{fig: stock time request}
مشاهده می کنید. در این روش به دلیل استفاده نکردن از تکنیک مقیاس پذیری خودکار،‌ زمان درخواست ها در سناریوهایی که تعداد درخواست های همزمان زیاد می باشد، مناسب نمی باشد.

\begin{figure}[t]
	\centering
	\includegraphics[scale=0.8]{stock-time-req.png}
	\caption{زمان استنتاج مدل در مسئله تحلیل احساسات دربازار سهام}
	\label{fig: stock time request}
\end{figure}








\subsection{مسئله توصیه گر کتاب}
\chapter{نتیجه‌گیری و پیشنهادات} \label{ch:Conclusion}

ما یک پلتفرم \lr{MLOps} متن‌باز و مبتنی بر ابر طراحی کرده‌ایم که خودکارسازی و بازتولید را در اکثر مراحل فرآیند یادگیری ماشین فراهم می‌کند. این پلتفرم در محیط‌های کوبرنتیز ارائه‌دهندگان مختلف ابری، کوبرنتیز داخلی و کوبرنتیز شبیه‌سازی شده روی ماشین‌های محلی اجرا می‌شود. هدف این طراحی، ایجاد زیرساختی است که مناسب اکثر پروژه‌ها و تیم‌های یادگیری ماشین باشد و فرآیندهای آن‌ها را به صورت خودکار و مقیاس‌پذیر نماید. 

محدودیت‌های فعلی و کارهای آینده شامل چندین بهبود مهم است که باید در پلتفرم \lr{MLOps} ما اعمال شود. پلتفرم برای اجرا به زیرساختی با منابع سخت‌افزاری زیاد نیاز دارد که توسعه محلی را دشوار و هزینه‌های سخت‌افزاری را افزایش می‌دهد؛ بنابراین، بهینه‌سازی‌های بیشتری برای سبک‌تر کردن پلتفرم ضروری است. بهینه‌سازی سرویس مش با جایگزینی \lr{Istio} با یک سرویس مش سبک‌تر، ادغام استفاده از سرویس مش با سرویس های کوبرنتیز به منظور افزایش سرعت پاسخ دهی به درخواست کاربر و کاهش اجزای غیرضروری \lr{Knative-Serving} نیز مورد نیاز است. امنیت سایبری پلتفرم باید تقویت شود و یک مطالعه موردی برای ارزیابی عملکرد و قابلیت اجرایی پلتفرم با پیاده‌سازی سیستم‌های یادگیری ماشین پیچیده و پرمنبع مانند آموزش مدل‌های زبانی بزرگ انجام شود. بهبودهای آینده شامل توسعه مرحله مدیریت شده برای برچسب‌گذاری داده‌ها، بهبودهای رابط کاربری برای ساده‌سازی استفاده از پلتفرم و ایجاد فرم‌های رابط کاربری برای افزودن استقرارها و جریان‌های کاری است.



% شروع مراجع %%%%%%%%%%%%%%%%%%%%%%%%%%%%%%%%%%%%%%%
% References
\begin{latin}
\renewcommand{\bibname}{\rl{مراجع}}
\bibliography{SepSrc}
% یکی از دو دستور زیر را فعال کنید. به نظر من برای یک تز یا یک کتاب باید مراجع به ترتیب الفبایی اسم نویسنده اول مرتب شوند، و مرتب کردن به ترتیب ارجاع در متن برای مقاله است.
%\bibliographystyle{IEEEsort} % مرتب کردن مراجع به ترتیب حروف الفبای اسم نویسنده اول
\bibliographystyle{IEEEbib}  % مرتب کردن مراجع به ترتیب ارجاع در متن
\end{latin}

% شروع قسمتهای انتهایی تز (چکیده انگلیسی و جلد انگلیسی) %%%%%%%%%%%%%%%%
\thispagestyle{empty}


\phantom{a}
\vfil

\begin{latin}
\centering
\begin{minipage}{0.9\textwidth}
{

%\noindent \rule{\textwidth}{2pt}

\noindent {\bf ABSTRACT}
\vskip 2mm

\baselineskip=0.4\baselineskip
In the current era, artificial intelligence and machine learning  have become vital and widely used technologies across various industries. These technologies enable companies and organizations to optimize processes, predict trends, and uncover hidden patterns in data with high accuracy and speed. However, fully leveraging the capabilities of AI and ML requires the effective and efficient deployment of ML models in production environments. MLOps, a combination of DevOps and ML concepts, aids in managing the lifecycle of ML models from development to deployment and maintenance. In this research, due to international sanctions and limited access to external services such as Google Vertex AI, Amazon SageMaker, and Microsoft Azure ML, a native MLOps platform has been designed and implemented. This platform is developed using open-source tools and provides various features including lifecycle management, performance monitoring, automatic updates, and integration with existing tools. The research aims to address questions related to the needs of a modern MLOps platform and the feasibility of its implementation using open-source tools. The results from the implementation and tests indicate that the proposed platform not only meets modern requirements but also ensures independence from cloud providers. Finally, the findings and recommendations for future research are presented.  \\
\baselineskip=2\baselineskip

\vspace{15mm}

\noindent {\bf KEYWORDS}
\vskip 2mm

{
1. Machine Learning Operation (MLOps). \\
2. Development and Operation (DevOps). \\
3. Continuous Integration (CI). \\
4. Continuous Development (CD). \\
5. Machine Learning Lifecycle Management. \\
6. Open-Source MLOps Platform. \\

}
}
\end{minipage}
\end{latin}
\vfil
  

\thispagestyle{empty}

\begin{latin}
\centering
\vspace*{-2.5cm}
\includegraphics[width=4cm]{Sharif_Logo.pdf}
\vspace*{-1.5cm}

{\large \bf SHARIF UNIVERSITY OF TECHNOLOGY} \\
{\bf ELECTRICAL ENGINEERING DEPARTMENT}

\vspace{2cm}

{\large \bf M.Sc. THESIS}
\vskip 1cm
{\sl Title:}\\

\baselineskip=2\baselineskip
{\huge \bf 

An Information-Theoretic Model for Image Watermarking
}\par

%\vspace{1.5cm}  % 1.2cm If both HAMKAAR and MOSHAVER exist
\vfil

by:\\
\vspace{-2mm}
{\large \bf AAAAA   BBBBBB}

%\vspace{2cm} % 1.1cm If both HAMKAAR and MOSHAVER exist
\vfil

{\sl Supervisor:}\\
\vspace{-1mm}
{\large \bf Dr. \dots\\}

%\vfil
%{\sl Co-supervisor:}\\
%\vspace{-1mm}
%{\large \bf  Dr. \dots\\}

%\vfil

%{\sl Advisor:}\\
%\vspace{-1mm}
%{\large \bf  Dr. \dots\\}

\vfil

{August 2005}
\par
\end{latin}



%%%%%%%%%%%%%%%%%%%%%%%%%%%%%%%%%%%%%%%%%%%%%%%%%%%%
\end{document}
